\begin{textsample}{POS Dim 2 – 2014\_c6 – Score 4.00 – 2014\_c6/enade\_2014\_mat\_c271511880.txt}  \label{ex:f2_pos_225}
CAPÍTULO 3 ANÁLISE TÉCNICA DA PROVA Este capítulo tem por objetivo apresentar o desempenho dos estudantes concluintes de Matemática ( Bacharelado e Licenciatura ) no ENADE/2014.
Para isso, foram calculadas as estatísticas básicas da prova em seu todo ( seção 3.1.1 ), bem como as estatísticas dos componentes relacionadas ao Componente de Formação Geral ( seção 3.1.2 ) e ao de Conhecimento Específico da Área ( seção 3.1.3 ).
Dadas as suas características, foram analisadas, em separado, as questões objetivas ( seção 3.2 ) e as questões discursivas ( seção 3.3 ) para concluintes de Bacharelado e de Licenciatura.
Para as questões objetivas, foram disponibilizados os índices de facilidade e o de discriminação bisserial, também em separado, para os Componentes de Formação Geral ( seção 3.2.1 ) e de Conhecimento Específico ( seção 3.2.2 ), sendo a análise da segunda seção separada por concluintes de Bacharelado e da Licenciatura.
De cada componente, uma das questões foi escolhida para exemplificar a análise gráfica, relacionando as alternativas escolhidas pelos estudantes ( inclusive o gabarito ) com o número de acertos no componente aferido, como já dito anteriormente, agregados por Bacharelado e Licenciatura.
O Anexo I apresenta a íntegra da análise gráfica para todas as questões objetivas agregadas e também separadas por Bacharelado e Licenciatura, a partir da questão 23.
Para cada uma das questões discursivas, os conteúdos dos tipos mais comuns de respostas dos estudantes são apresentados e comparados com o padrão de respostas esperado ( ver Anexo VIII com o padrão de respostas ).
Tomando como base as duas questões discursivas do Componente de Formação Geral, a seção 3.3.1.6 apresenta comentários sobre a correção das respostas com respeito à Língua Portuguesa.
Na Tabela 3.1 são apresentados o tamanho da população inscrita e de presentes, e as seguintes estatísticas das notas19 : média do desempenho na prova, erro padrão da média, desvio padrão, nota mínima, mediana e nota máxima.
Tais estatísticas contemplam o total de estudantes concluintes da Área de Matemática ( Bacharelado e Licenciatura ) inscritos e presentes à prova do ENADE/2014, tendo em vista as agregações por Grandes Regiões e o país como um todo.
As estatísticas das demais tabelas ímpares serão baseadas na mesma população da Tabela 3.1 e, portanto, do mesmo tamanho.
As tabelas pares são desagregadas por Categoria Administrativa e Organização Acadêmica da IES.
As estatísticas dessas tabelas pares serão baseadas na mesma população da Tabela 3.2. 19 Uma definição dessas estatísticas pode ser encontrada no Glossário.
Em relação aos gráficos de distribuição de notas, o intervalo considerado foi de 10 unidades, aberto à esquerda e fechado à direita, com exceção do primeiro intervalo, [ 0 ; 10 ], fechado em ambos os extremos.
Para os gráficos de distribuição das notas das questões discursivas, foram consideradas mais duas categorias : questão em branco e nota zero.
Todos os gráficos de distribuição de notas permitem a comparação dos resultados por habilitação : Bacharelado e Licenciatura. 3.1 ESTATÍSTICAS BÁSICAS DA PROVA Esta seção apresenta estatísticas selecionadas e histogramas da nota geral ( 3.1.1 ) e de cada componente : Formação Geral ( 3.1.2 ) e Conhecimento Específico ( 3.1.3 ).
São também apresentadas as estatísticas selecionadas de subpopulações caracterizadas por Grande Região, Categoria Administrativa e Organização Acadêmica. 3.1.1 Estatísticas Básicas Gerais A Tabela 3.1 apresenta as estatísticas básicas da prova, por grande Região, dos estudantes concluintes de Matemática ( Bacharelado e Licenciatura ).
A população total de inscritos foi de 17.347.
Destes, 13.854 estiveram presentes, sendo 20,1 \% o índice de não comparecimento.
A Região de maior abstenção foi a Norte ( 28,2 \% ), e a de menor abstenção foi a Centro-Oeste ( 18,0 \% ).
A média das notas da prova como um todo ( nas seções seguintes serão analisados os Componentes de Formação Geral e de Conhecimento Específico ) foi 32,5, sendo que os estudantes da região Norte obtiveram a média mais baixa ( 28,6 ), e os da região Sudeste obtiveram a média mais alta ( 34,4 ).
As demais médias foram : 31,3 na região Nordeste ; e 34,1 nas regiões Sul e Centro-Oeste.
O desvio padrão para o Brasil como um todo foi 11,6, sendo o maior desvio padrão encontrado na região Sudeste ( 12,4 ) e o menor, na região Norte ( 9,7 ), indicando uma dispersão um pouco menor das notas desta última região.
A região que obteve a maior nota máxima foi a Nordeste ( 90,4 ), ao passo que a região que atingiu a menor nota máxima foi a Norte ( 72,2 ).
A mediana do Brasil como um todo foi 31,4, sendo a maior mediana obtida na região Sudeste ( 33,2 ), e a menor, obtida na Norte ( 28,1 ).
A nota mínima foi zero em todas as regiões, sem exceção.
Considerando -se as notas segundo Grande Região, observa -se que existe diferença estatisticamente significativa ao nível de 95 \% entre a menor média, obtida na região Norte ( 28,6 ), e as médias de todas as demais regiões.
A Tabela 3.2 apresenta as estatísticas básicas da prova, desagregadas por Categoria Administrativa e por Organização Acadêmica.
Da população total de inscritos, 5.760 são de IES Privadas e 11.587 de IES Públicas.
Em relação à Organização Acadêmica, a maior participação foi obtida por estudantes de Universidades ( 13.163 ), seguido de Faculdades ( 2.329 ) e Centros Universitários ( 1.855 ).
A Categoria Administrativa de maior abstenção foi a Privada ( 20,8 \% ), e entre as Organizações Acadêmicas foi a dos Centros Universitários ( 25,2 \% ), ambas acima da média nacional de 20,1 \%.
A média das notas da prova como um todo foi 32,5.
Em relação à Categoria Administrativa, os estudantes das IES Públicas obtiveram média mais alta ( 32,8 ), e os das IES Privadas obtiveram média mais baixa ( 31,8 ), que a média nacional.
Observa -se que existe diferença estatisticamente significativa entre as médias das notas das IES Públicas e Privadas.
A diferença entre as médias das regiões Sudeste e Norte ( 5,8 ), a maior e a menor média regional, é superior à diferença entre IES Públicas e Privadas ( 1,0 ), caracterizando uma maior diversidade regional do que administrativa.
No tocante à Organização Acadêmica, apenas as Universidades obtiveram média mais alta ( 33,0 ) que a nacional.
As demais médias foram : 31,7 nos Centros Universitários e 30,1 nas Faculdades.
O desvio padrão para as IES Públicas ( 11,9 ) e para as Universidades ( também 11,9 ) foi superior ao do Brasil como um todo ( 11,6 ), indicando uma dispersão um pouco maior das notas nesta Categoria Administrativa e nesta Organização Acadêmica.
Constata -se que existe diferença estatisticamente significativa ao nível de 95 \% nas médias das notas dos estudantes provenientes de Universidades em relação aos de Centros Universitários e destes com os provenientes de Faculdades. 3.1.2 Estatísticas Básicas no Componente de Formação Geral A Tabela 3.3 apresenta as estatísticas básicas em relação ao componente da prova que avalia a Formação Geral dos estudantes concluintes.
Os estudantes de todo Brasil obtiveram desempenho médio de 51,8.
Quanto à variabilidade, o desvio padrão das notas dos estudantes do Brasil, como um todo, foi 18,3.
A maior média foi obtida na região Sul ( 53,7 ), e a menor, na região Norte ( 48,4 ).
As demais médias foram : 50,5 na região Nordeste ; 53,6 na região Sudeste ; e 52,7 na região Centro-Oeste.
Já o maior desvio padrão foi obtido nas regiões Nordeste e Sudeste ( 18,5 ), e o menor, na região Sul ( 17,6 ).
Os demais desvios padrões foram : 17,8 na região Norte ; e 18,4 na região Centro-Oeste.
A maior nota no Componente de Formação Geral da prova do ENADE/2014 foi 98,6, obtida por, pelo menos, um estudante na região Nordeste.
A menor nota máxima foi obtida nas regiões Norte e Centro-Oeste ( 97,8 ).
A mediana do Brasil, como um todo, foi 52,7, sendo a menor mediana encontrada na região Norte ( 49,2 ) ; e a maior encontrada na região Sudeste ( 55,2 ).
A nota mínima nesta parte foi zero em todas as regiões, sem exceção.
Considerando -se as notas segundo Grande Região, observa -se que existe diferença estatisticamente significativa entre a menor média das notas do Componente de Formação Geral, obtida na região Norte ( 48,4 ), e as médias de todas as demais regiões.
O Gráfico 3.2 propicia a avaliação do desempenho dos estudantes no Componente de Formação Geral a partir do histograma da distribuição das notas correspondentes ao Bacharelado e à Licenciatura.
As distribuições são unimodais, com a modas, respectivamente, nos intervalos ( 60 ; 70 ] e ( 50 ; 60 ].
Nota -se, ainda, que no Gráfico 3.2 as notas apresentam uma maior dispersão do que no Gráfico 3.1 ( distribuição das notas da prova ), confirmado pela comparação dos desvios padrões : 11,6 para a nota da prova como um todo e 18,3 para o Componente de Formação Geral.
Na Tabela 3.4 são apresentadas as informações referentes ao desempenho dos concluintes do Componente de Formação Geral, em diferentes agregações : Categoria Administrativa e Organização Acadêmica.
Observa -se que existe diferença estatisticamente significativa entre as médias dos tipos de Categoria Administrativa.
A maior média foi obtida por estudantes de IES Privadas ( 52,4 ), e a menor, obtida por estudantes de IES Públicas ( 51,5 ).
Considerando -se o tipo de Organização Acadêmica, nota -se que não há diferença estatisticamente significativa entre as médias dos estudantes de Centros Universitários ( 53,2 ) e de Universidades ( 52,0 ) ; mas existe diferença entre esses e os das Faculdades ( 49,7 ). 3.1.3 Estatísticas Básicas do Componente de Conhecimento Específico A Tabela 3.5 apresenta as estatísticas básicas referentes ao Componente de Conhecimento Específico da área de Matemática ( Bacharelado e Licenciatura ).
A média do desempenho dos estudantes do Brasil, como um todo, foi 26,0.
A maior média foi obtida na região Sudeste ( 28,0 ), e a menor, na região Norte ( 22,0 ).
As demais médias foram : 24,9 na região Nordeste ; 27,6 na região Sul ; e 27,9 na região Centro-Oeste.
Quanto à variabilidade das notas, o desvio padrão do Brasil, como um todo, foi 12,4, sendo o maior desvio padrão observado na região Sudeste ( 13,5 ), e o menor, na região Norte ( 9,9 ).
Os demais desvios foram : 11,4 na região Nordeste, 12,6 na região Sul e 13,3 na região Centro-Oeste.
A mediana das notas dos estudantes de todo o Brasil foi 23,9.
A maior mediana ocorreu nas regiões Sudeste, Sul e Centro-Oeste ( 26,9 em cada ), e a menor, na região Norte ( 22,4 ).
A mediana foi 23,4 na região Nordeste.
A nota máxima do Brasil, como um todo, foi 93,2, sendo obtida por, pelo menos, um estudante na região Sudeste.
As demais notas máximas foram : 74,7 na região Norte ; 89,8 na região Nordeste ; 85,3 na região Sul ; e 81,3 na região Centro-Oeste.
A nota mínima foi zero em todas as regiões, sem exceção.
Observa -se que existe diferença estatisticamente significativa da menor média das notas do Componente de Conhecimento Específico da região Norte ( 22,0 ) em relação às demais regiões.
A amplitude do intervalo de confiança para a média da região Centro-Oeste ( 1,8 ) é significativamente maior do que de todas as outras.
Assim como os Gráficos 3.1 e 3.2, o Gráfico 3.3, apresentado a seguir, permite uma avaliação do desempenho de concluintes dos cursos de Bacharelado e de Licenciatura em relação ao Componente de Conhecimento Específico com um histograma da distribuição das notas correspondentes.
Dentre as três distribuições apresentadas, esta é a mais concentrada nas notas baixas.
Estas também são distribuições unimodais, e o grupo modal é o ( 30 ; 40 ] para o Bacharelado, e o ( 20 ; 30 ] para a Licenciatura.
A Tabela 3.6 apresenta uma comparação dos resultados em relação à Categoria Administrativa e à Organização Acadêmica, agora levando em conta o desempenho de estudantes do Componente de Conhecimento Específico da prova.
No que se refere à Organização Acadêmica, a maior média foi das Universidades ( 26,7 ), vindo a seguir a dos Centros Universitários ( 24,5 ) e, depois, a das Faculdades ( 23,6 ).
O maior desvio padrão, e acima do valor para o Brasil, como um todo, foi o das Universidades ( 12,9 ).
As Universidades também obtiveram as maiores notas máxima ( 93,2 ) e mediana ( 24,4 ), seguidas, nessa ordem, pelos Centros Universitários ( 86,1 e 23,1 respectivamente ) e pelas Faculdades ( 72,1 e 22,4 respectivamente ).
A nota mínima foi zero para todas as Organizações Acadêmicas.
Quanto à Categoria Administrativa, observa -se um comportamento semelhante àquele da prova como um todo, ou seja, existe diferença estatisticamente significativa entre as médias das IES Públicas ( 26,6 ) e IES Privadas ( 24,9 ).
Neste caso, a maior média foi obtida por estudantes de IES Públicas de ensino.
Observa -se que não existem diferenças estatisticamente significativas ao nível de 95 \% no Componente de Conhecimento Específico entre as médias das notas dos Centros Universitários, e das Faculdades.
Existe, porém, diferença entre as médias destas categorias e das Universidades. 3.2 ANÁLISE DAS QUESTÕES OBJETIVAS Esta seção apresenta estatísticas para o conjunto de questões objetivas e uma análise gráfica de questões selecionadas dos componentes de Formação Geral ( 3.2.1 ), de Conhecimento Específico Bacharelado ( 3.2.2 ) e de Conhecimento Específico Licenciatura ( 3.2.3 ).
São também apresentadas e comparadas as médias das subpopulações caracterizadas por Grande Região. 3.2.1 Componente de Formação Geral A Tabela 3.7 apresenta as estatísticas básicas relativas às oito questões objetivas do componente da prova que abrange a Formação Geral dos estudantes.
A média do Brasil foi 53,4.
A menor média foi encontrada na região Norte ( 49,8 ), e a maior, na região Sul ( 55,3 ).
As demais médias foram : 52,4 na região Nordeste ; 54,9 na região Sudeste ; e 55,0 na região Centro-Oeste.
O desvio padrão do Brasil foi 21,6, sendo o maior desvio padrão encontrado na região Nordeste ( 21,8 ), e o menor, na região Sul ( 20,7 ).
Os demais desvios foram : 21,4 na região Norte ; e 21,6 nas regiões Sudeste e Centro-Oeste.
As medianas ( 50,0 ), as notas máximas ( 100,0 ) e as notas mínimas ( 0,0 ) foram iguais para todas as regiões, exceto a mediana nas regiões Sudeste e Sul ( 62,5 em ambas ).
A Tabela 3.8 apresenta o índice de facilidade e o índice de discriminação ( ponto bisserial ) para cada uma das questões objetivas do Componente de Formação Geral.
Quanto ao índice de facilidade, foram usadas as seguintes cores para diferenciar o nível de dificuldade da questão :  Azul para as questões classificadas com índice muito fácil ( \&gt ; =0,86 ), verde para as questões classificadas com índice fácil ( 0,61 a 0,85 ), amarelo para as questões classificadas com médio ( 0,41 a 0,60 ), vermelho para as questões classificadas com difícil ( 0,16 a 0,40 ) e roxo para as questões classificadas com muito difícil ( \&lt ; =0,15 ).
Já quanto ao índice de discriminação, foram usadas as seguintes cores para qualificar a questão :  As questões classificadas com índice fraco receberam a cor vermelho ( \&lt ; =0,19 ), as classificadas com médio receberam a cor amarelo ( 0,20 a 0,29 ), as classificadas com bom receberam a cor verde ( 0,30 a 0,39 ) e as classificadas com muito bom ( \&gt ; =0,40 ) receberam a cor azul.
As questões objetivas do Componente de Formação Geral, segundo o índice de facilidade, foram assim avaliadas : das oito questões, nenhuma teve o índice de facilidade classificado como muito fácil.
Duas questões foram tidas como fácil, por terem índice de acertos situado na faixa entre 0,61 e 0,85 ( de 61,0 \% a 85,0 \% de acertos ).
Cinco questões foram consideradas de dificuldade média, situando -se no intervalo entre 0,41 e 0,60 do índice de facilidade, ou seja, houve entre 41,0 \% e 60,0 \% de acertos, enquanto uma questão foi classificada na categoria difícil, situando -se no intervalo entre 0,16 e 0,40.
Por fim, nenhuma das questões apresentou menos de 15 \% de acertos, razão pela qual seria classificada como muito difícil.
Como já comentado, para análise das questões objetivas, relativas à Formação Geral, segundo o poder de discriminação, utilizou -se o índice de discriminação ( ponto bisserial ).
Nesta análise, para esse grupo de estudantes, as questões foram assim avaliadas : cinco das oito questões apresentaram índice acima ou igual a 0,40 e, assim, foram classificadas com índice muito bom ; três questões tiveram índice de discriminação bom, com valor entre 0,30 e 0,39 ; e nenhuma questão teve nível médio ou fraco de discriminação.
O índice de facilidade variou de 0,31 a 0,77, e o de discriminação, de 0,37 a 0,53.
As cinco questões com índice de discriminação muito bom figuraram entre os diversos níveis de dificuldade desse conjunto : duas classificadas na categoria fácil ( questões 4 e 8 ) do índice de facilidade, e três, na categoria médio ( questões 3, 6 e 7 ).
Em particular, a questão 7 foi a que apresentou o maior poder discriminatório, com índice 0,53, e foi considerada média em termos de facilidade, com uma proporção de 0,47 acertos.
O máximo de acertos foi alcançado pela questão 8 com um índice de facilidade de 0,77.
A questão de número 1 apresentou um índice de facilidade de 0,48, ou seja, quase a metade dos estudantes conseguiu resolvê -la, dentro do universo de participantes.
Seu índice de discriminação foi bom ( 0,39 ).
Já a questão 5 também obteve índice de discriminação bom, 0,38, porém seu índice de facilidade foi difícil ( 0,31 ).
O Gráfico 3.4, para exemplificar, analisa o comportamento da questão de número 7 de Formação Geral.
Trata -se de uma questão considerada de facilidade média e a que obteve o maior índice de discriminação dessa parte da prova.
Neste gráfico, cada uma das cinco curvas representa o percentual de respostas em determinada alternativa da questão, em função do número de acertos dos estudantes nessa parte da prova ( Formação Geral/Múltipla Escolha ), antes de possíveis eliminações pelo critério do ponto bisserial.
A curva em vermelho corresponde à alternativa E, a correta para esta questão.
Assim, observa -se que entre os estudantes com menor número de acertos, nessa parte do exame, a situação mais frequente foi a escolha de uma das alternativas incorretas : a alternativa C ( em preto ) ou D ( em roxo ).
À medida em que o número de acertos aumenta, indicando desempenho melhor nessa parte da prova, aumenta concomitantemente a proporção de estudantes que selecionaram a alternativa correta E, atingindo 100 \% somente para os estudantes com 8 acertos.
Essa análise permite verificar como a questão discriminou os grupos de desempenho, justificando o alto índice obtido na questão.
Os gráficos relativos às demais questões de Formação Geral constam do Anexo I. 3.2.2 Componente de Conhecimento Específico – Matemática ( Bacharelado ) A Tabela 3.9 apresenta as estatísticas básicas em relação às questões objetivas do Componente de Conhecimento Específico da prova de Matemática ( Bacharelado ), por Grande Região.
A média do Brasil deste componente foi de 40,2.
A menor média foi observada na região Sul ( 35,6 ), e a maior, na região Centro-Oeste ( 46,5 ).
O desvio padrão de todo o Brasil foi 19,5, sendo o menor desvio padrão encontrado na região Centro-Oeste ( 17,1 ), e o maior, na região Norte ( 20,8 ).
A mediana de todo o Brasil foi 40,0, a mesma encontrada nas regiões Nordeste, e Sudeste.
As regiões Norte e Sul ( 32,0 em ambas ) apresentaram valores menores para a mediana, enquanto a região Centro-Oeste ( 44,0 ) apresentou valor maior.
A nota máxima da prova foi 92,0, obtida nas questões objetivas do Componente de Conhecimento Específico, por, pelo menos, um estudante na região Sudeste, enquanto nas demais regiões, a nota máxima da prova foi : 72,0 na região Norte ; 88,0 na região Nordeste ; e 80,0 nas regiões Sul e Centro-Oeste.
Nas regiões Sudeste e Sul a nota mínima foi zero, enquanto nas demais regiões foi : 16,0 na região Norte ; 12,0 na região Nordeste ; e 8,0 na região Centro-Oeste.
A Tabela 3.10 apresenta os índices de facilidade e discriminação ( ponto bisserial ) das questões objetivas do Componente de Conhecimento Específico para os estudantes de Matemática ( Bacharelado ).
Para facilitar a diferenciação das questões usaram -se as mesmas cores da Tabela 3.8 para as diferentes classificações dos índices de facilidade e de discriminação.
Dentre as questões objetivas da parte da prova relativa ao Componente de Conhecimento Específico, nenhuma das questões foi anulada pela Comissão.
Desse modo, a classificação, quanto ao índice de facilidade, foi estabelecida com base em todas as 27 questões.
A partir dos índices obtidos, pode -se concluir que mais de um terço das questões objetivas da prova foi considerado, pelo menos, difícil : das 27 questões, 15 foram classificadas como difícil ( a classificação modal ) e uma como muito difícil.
Duas questões foram classificadas como fácil, e outras nove consideradas como médio.
Nenhuma foi classificada como muito fácil Já quanto aos índices de discriminação das questões objetivas do Componente de Conhecimento Específico da prova, tem -se como resultado a seguinte classificação : oito das 27 questões foram consideradas como boas, enquanto 15 delas tiveram índice de discriminação muito bom.
Assim, para a maioria das questões – 23 em 27 – os índices de discriminação foram bom ou muito bom.
Dentre as demais, duas delas foram classificadas como médio e outras duas como fraco, sendo quatro, por conseguinte, a quantidade de questões nos dois patamares mais baixos de discriminação.
Constata -se, assim, que a prova – no que se refere ao Componente de Conhecimento Específico – possuía muito boa capacidade de discriminar entre aqueles que dominam ou não o conteúdo.
Dentre as questões que alcançaram os maiores índices de discriminação, 15 delas, as de números 10, 11, 13, 15, 16, 17, 18, 22, 23, 24, 25, 26, 28, 32 e 35, foram classificadas com índice de discriminação muito bom, situando -se no intervalo de 0,40 a 0,55 do índice, duas delas ( questões 17 e 18 ) foram classificadas na categoria fácil, quanto ao índice de facilidade, sete delas ( questões 10, 13, 15, 23, 24, 28 e 35 ), na categoria médio, e as outras seis ( questões 11, 16, 22, 25, 26 e 32 ), na categoria difícil.
A questão de número 20 foi a mais difícil dentre as 27 questões específicas, com baixo índice de facilidade, apenas 12,0 \% de acertos.
Essa questão apresentou poder discriminatório igualmente baixo, 0,06, o que comprova ter sido esta questão a mais difícil para os estudantes.
Destaca -se, também, a questão 21, com índice de facilidade 0,16, o que, em termos percentuais, corresponde a 16,0 \% de estudantes que responderam acertadamente e que também teve 0,06 de índice de discriminação.
Tais questões foram, portanto, pelo critério ponto bisserial, consideradas inadequadas.
Por isso, as questões 20 e 21 foram as únicas eliminadas do cômputo da nota final.
A título de exemplo das análises do comportamento das questões objetivas, o Gráfico 3.5 analisa a questão 28 do Componente de Conhecimento Específico.
Esta questão foi classificada como de facilidade média.
Apresentou índice de facilidade 0,50, ou seja, 50,0 \% dos estudantes assinalaram acertadamente a opção E, correspondente ao gabarito.
Seu índice de discriminação foi igual a 0,55, classificado como muito bom, também sendo esta questão a que apresentou o maior índice discriminatório.
Neste gráfico, cada uma das cinco curvas representa o percentual de respostas em determinada alternativa da questão 28, em função do número de acertos dos estudantes nessa parte da prova, antes de possíveis eliminações de questões pelo critério do ponto bisserial.
A alternativa correta E, representada no gráfico pela curva em vermelho, foi escolhida em maiores proporções pelos estudantes com desempenho melhor nessa parte da prova.
Já as alternativas incorretas, também denominadas distratores, foram selecionadas, principalmente, por aqueles com notas mais baixas.
Neste caso a soma não é sempre 100 \% por causa das questões não respondidas ou com mais de uma opção marcada.
Aqueles com nota zero, na sua quase totalidade deixaram esta questão em branco ou marcaram mais de uma alternativa, comportamento considerado inválido.
A proporção de estudantes que selecionou a resposta correta E aumenta gradativamente, chegando a atingir 100 \% para 19 acertos e para 21 acertos ou mais, enquanto a proporção dos que escolheram alternativas incorretas decai, a partir de três acertos, como função do número de acertos nessa parte da prova.
Os gráficos relativos às demais questões do Componente de Conhecimento Específico constam do Anexo I. 3.2.3 Componente de Conhecimento Específico – Matemática ( Licenciatura ) A Tabela 3.11 apresenta as estatísticas básicas em relação às questões objetivas do Componente de Conhecimento Específico da prova de Matemática ( Licenciatura ), por Grande Região.
A média do Brasil deste componente foi de 28,4.
A menor média foi observada na região Norte ( 24,8 ), e a maior, na região Sudeste ( 30,1 ).
O desvio padrão de todo o Brasil foi 12,7, sendo o menor desvio padrão encontrado na região Norte ( 10,9 ), e o maior, na região Centro-Oeste ( 13,6 ).
A mediana de todo o Brasil foi 26,3, a mesma encontrada em todas as regiões.
A nota máxima da prova foi 89,5, obtida nas questões objetivas do Componente de Conhecimento Específico, por, pelo menos, um estudante na região Nordeste, enquanto nas demais regiões, a nota máxima da prova foi : 73,7 na região Norte ; e 84,2 nas regiões Sudeste, Sul e Centro-Oeste.
Em todas as regiões a nota mínima foi zero, sem exceção.
A questão de A Tabela 3.12 apresenta os índices de facilidade e discriminação ( ponto bisserial ) das questões objetivas do Componente de Conhecimento Específico para os estudantes de Matemática ( Licenciatura ).
Para facilitar a diferenciação das questões usaram -se as mesmas cores das Tabelas 3.8 e 3.10 para as diferentes classificações dos índices de facilidade e de discriminação.
Dentre as questões objetivas da parte da prova relativa ao Componente de Conhecimento Específico, nenhuma das questões foi anulada pela Comissão.
Desse modo, a classificação, quanto ao índice de facilidade, foi estabelecida com base em todas as 27 questões.
A partir dos índices obtidos, pode -se concluir que quase todas as questões objetivas da prova da Licenciatura foram consideradas, pelo menos, difíceis : das 27 questões, 25 foram classificadas como difícil ( a classificação modal ) ou como muito difícil.
Duas questões foram classificadas como médio.
Já quanto aos índices de discriminação das questões objetivas do Componente de Conhecimento Específico da prova, tem -se como resultado a seguinte classificação : 5 das 27 questões foram consideradas como boas, enquanto nenhuma delas teve índice de discriminação muito bom.
Assim, para a minoria das questões – 5 em 27 – os índices de discriminação foram bom.
Dentre as demais, 14 delas foram classificadas como médio e outras oito como fraco, sendo 22, por conseguinte, a quantidade de questões nos dois patamares mais baixos de discriminação.
Constata -se, assim, que a prova – no que se refere ao Componente de Conhecimento Específico de Licenciatura – possuía capacidade fraca de discriminar entre aqueles que dominam ou não o conteúdo.
Dentre as questões que alcançaram os maiores índices de discriminação, cinco delas, as de números 17, 18, 25, 31 e 32, foram classificadas com índice de discriminação bom, situando -se no intervalo de 0,30 a 0,39 do índice, duas delas ( questões 18 e 31 ) foram classificadas na categoria médio, quanto ao índice de facilidade, e as outras três ( questões 17, 25 e 32 ), na categoria difícil.
A questão de número 16 foi a mais difícil dentre as 27 questões específicas, com baixo índice de facilidade, apenas 9,0 \% de acertos.
Essa questão apresentou poder discriminatório igualmente baixo, 0,14, o que comprova ter sido esta questão a mais difícil para os estudantes.
Destaca -se, também, a questão 27, com índice de facilidade 0,15, o que, em termos percentuais, corresponde a 15,0 \% de estudantes que responderam acertadamente.
Já o índice de discriminação foi 0,12.
Tais questões foram, portanto, pelo critério ponto bisserial, consideradas inadequadas.
Por isso, as questões 16 e 27 foram eliminadas do cômputo da nota final.
Além destas duas, as demais questões com índice de discriminação fraco, questões 20, 21, 22, 28, 29 e 33 também não foram utilizadas no cômputo final das notas, num total de oito questões eliminadas.
A título de exemplo das análises do comportamento das questões objetivas, o Gráfico 3.6 analisa a questão 18 do Componente de Conhecimento Específico.
Esta questão apresentou índice de facilidade 0,53, ou seja, 53,0 \% dos estudantes assinalaram acertadamente a opção A, correspondente ao gabarito.
Seu índice de discriminação foi igual a 0,33, classificado como bom, também sendo esta questão a que apresentou o maior índice discriminatório.
Neste gráfico, cada uma das cinco curvas representa o percentual de respostas em determinada alternativa da questão 18, em função do número de acertos dos estudantes nessa parte da prova, antes de possíveis eliminações de questões pelo critério do ponto bisserial.
A alternativa correta A, representada no gráfico pela curva em azul, foi escolhida em maiores proporções pelos estudantes com desempenho melhor nessa parte da prova.
Já as alternativas incorretas, também denominadas distratores, foram selecionadas, principalmente, por aqueles com notas mais baixas.
Neste caso também a soma não é sempre 100 \% por causa das questões não respondidas ou com mais de uma opção marcada.
Aqueles com nota zero, na sua quase totalidade deixaram esta questão em branco ou marcaram mais de uma alternativa, comportamento considerado inválido.
A proporção de estudantes que selecionou a resposta correta A aumenta gradativamente, chegando a atingir 100 \% para 17 acertos ou mais, enquanto a proporção dos que escolheram alternativas incorretas decai, a partir do primeiro acerto, como função do número de acertos nessa parte da prova.
Os gráficos relativos às demais questões do Componente de Conhecimento Específico constam do Anexo I. 3.3 ANÁLISE DAS QUESTÕES DISCURSIVAS Esta seção apresenta estatísticas selecionadas e histogramas das Questões Discursivas de Formação Geral ( 3.3.1 ) e Conhecimento Específico ( 3.3.2 ).
São também apresentadas e comparadas as médias de subpopulações caracterizadas por Grande Região. 3.3.1 Componente de Formação Geral As análises dos resultados de desempenho dos estudantes de Matemática ( Bacharelado e Licenciatura ) nas duas questões discursivas relativas à Formação Geral, considerando -se as notas de conteúdo e de desempenho linguístico, encontram -se na Tabela 3.13 e no Gráfico 3.7.
Na Tabela 3.13, observa -se que a nota média nesse conjunto de questões foi próxima da obtida nas objetivas.
Os estudantes, de todo o Brasil, obtiveram, em Formação Geral, média 53,4 nas questões objetivas e 49,3 nas questões discursivas.
No entanto, pode -se notar um aumento do desvio padrão de 21,6 nas questões objetivas do Componente de Formação Geral dos estudantes de todo o Brasil, para 25,7 nas questões discursivas do mesmo componente.
A maior média foi obtida na região Sudeste ( 51,7 ), e a menor, na região Norte ( 46,2 ).
A mediana de todo o Brasil, neste componente, foi 54,5.
Nas regiões Sudeste e Sul a mediana foi maior, 57,0 e 56,0, respectivamente.
Nas demais regiões foi menor : Norte ( 50,0 ), Nordeste ( 52,5 ) e Centro-Oeste ( 54,5 ).
A nota máxima ( 99,0 ) foi obtida por pelo menos um estudante na região Norte, sendo 97,5 nas regiões Nordeste e Sudeste, 98,0 na região Sul e 97,0 na região Centro-Oeste.
A nota mínima ( 0,0 ) foi a mesma em todas as regiões do Brasil, sem exceção.
O Gráfico 3.7 representa a distribuição das notas nas questões discursivas do Componente de Formação Geral, segundo a opção de Bacharelado/Licenciatura.
As modas destas distribuições ocorrem no intervalo [ 0 ; 10 ] para o Bacharelado e no intervalo ( 60 ; 70 ] para Licenciatura.
Na sequência, os resultados verificados para cada uma das questões discursivas de Formação Geral serão apresentados, estabelecendo -se relações com os conteúdos abordados em cada uma delas.
Os comentários da Banca de docentes corretores, a respeito do observado na correção das respostas dos estudantes, suas impressões e conclusões serão apresentados junto à análise de cada questão.
Cumpre esclarecer que, tendo em vista que as questões discursivas de Formação Geral são padronizadas, ou seja, constam de todas as provas, os comentários da Banca são os mesmos para todas as carreiras acadêmicas, sendo direcionados a todos os estudantes que participaram do ENADE/2014.
A seguir, serão analisados os desempenhos da Área de Matemática ( Bacharelado e Licenciatura ) nas duas questões discursivas de Formação Geral do ENADE/2014, comparando -se os resultados obtidos com comentários para cada questão. 3.3.1.1 Análise de conteúdo da Questão Discursiva 1 do Componente de Formação Geral Os dados de Matemática ( Bacharelado e Licenciatura ), obtidos a partir das respostas à questão 1, encontram -se na Tabela 3.14 e no Gráfico 3.8.
Nessa questão – de mesmo desempenho médio que a questão discursiva 2 de Formação Geral – os estudantes, de todo o Brasil, tiveram média 48,5.
A maior média para a questão 1 foi obtida na região Sudeste ( 52,4 ), e a menor, na região Norte ( 42,4 ).
Quanto à variabilidade das notas, o desvio padrão de todo o Brasil foi 32,1.
O menor desvio padrão foi obtido na região Norte ( 30,8 ), e o maior desvio padrão foi obtido na região Centro-Oeste ( 32,5 ).
As medianas das regiões Nordeste, Centro-Oeste e do Brasil, como um todo, foram iguais ( 50,0 ).
Nas regiões Sudeste e Sul, a mediana foi 55,0 ; e na região Norte, foi 45,0.
As notas máximas e mínimas da questão discursiva 1 foram as mesmas para todas as regiões do Brasil, respectivamente, 100,0 e 0,0.
O Gráfico 3.8 mostra a distribuição das notas na questão discursiva 1 do Componente de Formação Geral para estudantes do Bacharelado e da Licenciatura.
Observa -se que a maior frequência corresponde aos estudantes que deixaram a questão em branco, para o Bacharelado, e aos que obtiveram notas no intervalo ( 70 ; 80 ], para a Licenciatura. 3.3.1.2 Comentários sobre a correção de Conteúdo das respostas à Questão Discursiva 1 O enunciado apresentado da Questão 1 era claro e foi redigido em linguagem acessível.
A questão abordava o tema da mobilidade urbana, relacionando -a aos ideais de desenvolvimento urbano sustentável estabelecidos pela ONU, do ponto de vista ecológico, cultural, político, institucional, social e econômico.
O enunciado apresentava uma tabela, sem referência à fonte de publicação, que detalhava percentualmente as modalidades de deslocamento da população urbana brasileira, com base em duas distinções : motorizado versus não motorizado ; coletivo versus individual.
O comando da questão solicitava, então, que o estudante desenvolvesse um texto dissertativo composto de duas partes : análise das consequências do transporte motorizado para o desenvolvimento sustentável e elaboração de duas propostas de intervenção pautadas pelo incentivo ao uso da bicicleta, meio de transporte não motorizado e, portanto, não poluente.
A temática dos efeitos da industrialização, com foco na produção de automóveis, e consequências para o meio ambiente é atual.
Além disso, a questão provoca a reflexão sobre a necessidade de adequação das cidades ao número de habitantes, no contexto da utilização dos meios de transportes de massa : metroviário, ferroviário, por ônibus.
Esperava -se, assim, que o concluinte recorresse a diferentes argumentos para analisar as consequências do transporte motorizado, explicitados pelo texto motivador, e propusesse duas ações de intervenção, as quais considerasse relevantes para o incremento ao uso de bicicleta, que, preferencialmente, deveriam manter coerência com a argumentação desenvolvida na primeira parte da resposta.
O conflito entre a busca pelo transporte individual e o coletivo é inerente a todas as grandes cidades do mundo.
No Brasil, esse conflito emerge com maior gravidade, como efeito do aumento da quantidade de veículos automotores ( automóveis e motocicletas ) e da baixa qualidade dos transportes coletivos.
Tal associação incentiva o uso de meios de transporte individual em detrimento da utilização dos coletivos.
O padrão de respostas era tecnicamente adequado ao enunciado e amplo o suficiente para abranger as variações de respostas apresentadas.
Deve ser ressaltado que a proteção ao ambiente ecologicamente equilibrado é tema conhecido dos concluintes dos cursos universitários, que deveriam possuir informações suficientes para desenvolver uma resposta abordando pelo menos dois dos aspectos listados no padrão de respostas como possibilidades de consequências ( item ‘ a ’ ) e duas propostas de ação de intervenção aceitáveis ( item ‘ b ’ ).
Segundo o padrão, concluintes poderiam desenvolver, quanto ao item ‘ a ’, os seguintes pontos : aumento da emissão de poluentes atmosféricos ; aumento da emissão de gases de efeito estufa ( CO2- dióxido de carbono, CO - monóxido de carbono, O3 - Ozônio ) ; aumento da poluição visual e sonora ; aumento da temperatura local e global ; aumento do consumo de combustíveis ; aumento de problemas de saúde ( cardíacos, respiratórios, dermatológicos ) ; aumento da frota de veículos, promovendo congestionamentos urbanos ; diminuição de áreas verdes ; desmatamento ; aumento das áreas de impermeabilização, resultando em enchentes, diminuição da infiltração da água e recarga de lençóis freáticos ; elevação dos custos de manutenção das cidades ( metroferrovias, rodovias, tratamento de água, limpeza da cidade, etc. ) ; necessidade de ampliação de vias trafegáveis ; necessidade de ampliação de áreas de estacionamento.
Em relação ao item ‘ b ’, os estudantes poderiam considerar os seguintes aspectos : construção de vias exclusivas para bicicletas ( ciclovias e ciclofaixas ) ; proposição de formas de integração entre o transporte por bicicletas, o metroviário e os ônibus coletivos, a fim de garantir segurança e conforto em momentos de adversidades climáticas e relevo acidentado ; pontos de aluguel e/ou empréstimos de bicicletas ; construção de bicicletários ; investimentos na segurança pública ; políticas de incentivo ao uso de bicicletas ( educação ambiental, qualidade de vida, saúde, propaganda ) ; implementação de políticas de crédito e de redução de custos das bicicletas.
Para o item ‘ a ’, foi encontrado um número expressivo de respostas adequadas.
As consequências mais citadas estavam relacionadas com a emissão de poluentes e o aumento dos engarrafamentos.
Houve críticas frequentes à qualidade do transporte público e à necessidade de locomoção rápida diante da distância das residências aos locais de trabalho.
Em tom de reivindicação, muitos estudantes defenderam solução para os problemas do transporte coletivo, quer metroviário, ferroviário ou por ônibus.
O uso da bicicleta também foi considerado por muitos como um meio de transporte benéfico à saúde, mas alguns apontaram a impossibilidade de sua utilização pelos idosos e deficientes.
No concernente ao item ‘ b ’, de modo quase uniforme citou -se a necessidade de ciclovias e ciclofaixas, bem como de se garantir segurança para os ciclistas.
A necessidade de manutenção das ciclovias e ciclofaixas existentes em algumas cidades foram também lembradas.
Com menos frequência também houve a indicação de maior conscientização da população, inclusive pela mídia, e pelo incentivo à aquisição desses veículos não motorizados, por meio da diminuição de impostos para redução do custo de aquisição.
Houve quem sugerisse, inclusive, uma política de doação de bicicletas.
Alguns mencionaram a necessidade de se assegurar integração das ciclovias com os demais modais, com a construção de bicicletários seguros diante das grandes distâncias a percorrer, notadamente nas metrópoles.
Também com base na dificuldade de o veículo não motorizado vencer as longas distâncias que o trabalhador deve percorrer nos centros urbanos para chegar ao local onde exerce a sua atividade produtiva, alguns estudantes argumentaram contra a utilização da bicicleta como forma de resolução da poluição ambiental.
O principal equívoco cometido na sugestão de ações ( item ‘ b ’ ) foi o de mencionar a necessidade de melhoria do meio ambiente dissociado do uso de bicicletas, como se requeria no comando da questão.
Nesse sentido, foram encontradas respostas que indicavam a caminhada como opção.
Foram consideradas respostas fracas ( notas de zero a 30 ) aquelas meramente opinativas, ou que apenas transcreveram elementos do enunciado, sem qualquer acréscimo ou apreciação crítica.
As respostas medianas ( notas entre 35 e 70 ) foram as que não apresentavam algum desenvolvimento para os tópicos citados ( duas consequências e duas ações de intervenção ) ou por falta de alguns deles.
Já as respostas boas ( notas de 75 a 100 ) continham os tópicos considerados pelo padrão de respostas, com algum desenvolvimento, pelo menos, regular de argumentação.
As respostas dos estudantes revelaram um adequado conhecimento quanto aos efeitos do transporte motorizado para o ambiente e que a temática da ecologia parece ser tema bem compreendido, aparecendo nas respostas como uma preocupação para as futuras gerações.
Nesse sentido, algumas políticas públicas foram mencionadas apropriadamente para assegurar o desenvolvimento sustentável pelos respondentes, tais como : educação, habitação, saneamento, saúde, transporte e segurança.
No entanto, de modo geral, as respostas demostraram que os estudantes têm muita dificuldade na expressão escrita do pensamento, como se pode constatar pela avaliação do desempenho linguístico que ficou a cargo de uma banca específica, formada por profissionais da área de Língua Portuguesa. 3.3.1.3 Análise de Conteúdo da Questão Discursiva 2 do Componente de Formação Geral A Tabela 3.15 mostra que o desempenho médio dos estudantes na questão discursiva 2 ( média 48,5 ) foi o mesmo obtido na questão discursiva 1 ( média 48,5 ).
A região Sudeste foi também nessa questão, aquela cuja média foi maior ( 49,8 ).
As de menor média foram as regiões Norte, Nordeste e Centro-Oeste ( 47,6 ).
Quanto à variabilidade das notas, o desvio padrão de todo o Brasil foi 32,4, pouco superior ao obtido na questão discursiva 1 ( 32,1 ).
O maior desvio nessa questão foi obtido na região Norte ( 32,8 ), enquanto o menor foi obtido na região Sul ( 31,9 ).
As medianas ( 50,0 ), as notas máximas ( 100,0 ) e as notas mínimas ( 0,0 ) foram iguais para todas as regiões, exceto a mediana na região Sudeste ( 60,0 ).
O Gráfico 3.9 mostra a distribuição das notas na questão discursiva 2 do Componente de Formação Geral para estudantes do Bacharelado e da Licenciatura.
Observa -se que a maior frequência corresponde aos estudantes que deixaram a questão em branco tanto para aqueles de Bacharelado quanto para os de Licenciatura.
Nota -se, ainda, que as notas ficaram mais distribuídas ao longo dos intervalos em comparação à questão discursiva de número 1. 3.3.1.4 Comentários sobre a correção de Conteúdo das respostas à Questão Discursiva 2 A questão 2 abordava o tema da violência urbana a partir de um texto motivador recolhido em um site de um jornal brasileiro.
O fragmento relatava a prisão de três jovens por agressão a um outro jovem, menor de idade, guardador de carros.
A notícia trazia informações sobre o estado do jovem agredido e sobre as possíveis penas a serem aplicadas aos agressores.
O enunciado solicitava, então, que o estudante desenvolvesse um texto dissertativo composto de duas partes : análise de duas causas do tipo de violência descrito no texto e apresentação de dois fatores que contribuiriam para evitar em situações como a relatada.
Como o texto motivador reduzia -se ao relato de um episódio de violência urbana, o enunciado não oferecia argumentos que ajudassem o estudante a elaborar uma análise das possíveis causas do fato.
As propostas a serem apresentadas, no sentido de se evitar em situações como a relatada, deveriam manter coerência com a análise das causas, desenvolvida na primeira parte da resposta.
O enunciado era claro e a matéria permitia identificar tratar -se de um conflito tipicamente urbano no qual se envolveram jovens de baixa renda em disputa por um território onde exerciam trabalho informal.
Sendo essa uma das hipóteses albergada sobre o amplo tema da violência urbana no qual se desenvolve o relato.
A temática da questão é bastante conhecida, não somente por debates no âmbito universitário, especialmente na área de Ciências Sociais, bem como pela ênfase que é dada pela mídia onde assume contornos, muitas vezes, maiores do que o problema em si.
Assim, por tratar -se de tema bastante presente e debatido no cotidiano das pessoas, a questão pode ser considerada de baixa complexidade.
O padrão de resposta apresentava uma gama de possibilidades de causas da violência aceitas como corretas para o item a ’, tais como : problemas relacionados à educação ( baixa escolaridade, evasão escolar, qualidade da educação, distanciamento entre escola e a realidade social, tempo de permanência na escola ) ; desigualdades socioculturais ( gênero, etnia, economia, etc. ) ; desemprego e falta de qualificação profissional ; precariedade da segurança pública ; uso de drogas ; desvalorização da vida humana ; banalização da violência ; sensação de impunidade ; ausência de políticas sociais ; degradação da vida urbana ; desconhecimento ou desrespeito aos direitos humanos e constitucionais ; desestruturação familiar ; desvalorização de princípios éticos e morais.
No concernente ao item ‘ b ’, o padrão também previa a possibilidade de apresentação de diversos fatores que podem contribuir para evitar a violência : políticas de segurança mais efetivas ; políticas públicas de melhoria das condições socioeconômicas ; maior consciência cidadã e respeito à vida ; melhor distribuição de renda ; melhoria da educação ( aumento da escolaridade, redução da evasão escolar, qualidade da educação, aproximação entre a escola e a realidade social, aumento do tempo de permanência na escola ) ; aumento da oferta de emprego e melhoria da qualificação profissional ; medidas preventivas ao uso de drogas ; maior eficácia do sistema judiciário ; revisão da legislação penal ; valorização de princípios éticos, morais e familiares.
O enunciado da questão permitia que os estudantes, em suas respostas, pautassem o problema por diversos ângulos.
Dessa forma, o padrão buscou indicar diversas possibilidades a serem identificadas pelo corretor na leitura da resposta do estudante.
As notas foram atribuídas considerando o desenvolvimento da resposta e a indicação de ao menos dois dos itens exigidos pelo comando da questão nos itens ’ a ’ ( duas causas ) e ‘ b ’ ( dois fatores ), dentre os considerados pelo padrão de respostas.
No entanto, o padrão exigia que as respostas estivessem afinadas com o perfil profissional previsto no Art. 3º da Portaria nº 255, de 02 de junho de 2014, onde se lê : Art. 3º No componente de Formação Geral serão considerados os seguintes elementos integrantes do perfil profissional : atitude ética ; comprometimento social ; compreensão de temas que transcendam ao ambiente próprio de sua formação, relevantes para a realidade social ; espírito científico, humanístico e reflexivo ; capacidade de análise crítica e integradora da realidade ; e aptidão para socializar conhecimentos em vários contextos e públicos diferenciados.
Nesse sentido, teses repressivas não conformadas com o texto constitucional e amparados pelos direitos humanos, não foram consideradas.
Questões como a majoração das penas, instituição de penas mais duras e criminalização da conduta dos menores desviados, só foram consideradas no caso de serem tratadas como menção a um debate que está presente, nunca como medida a ser aplicada de forma arbitrária e inconstitucional.
Também não foram previstas visões mais autoritárias e repressoras como a defesa de proibição do trabalho informal, como é o caso dos “ flanelinhas ”, e a defesa de não fornecer dinheiro aos “ guardadores ” de carro, identificado como esmolas por alguns.
Não foram consideradas respostas que meramente reproduziam o enunciado sem qualquer reflexão ou análise.
Aquelas que utilizaram os elementos do enunciado, mesmo que apenas transcritos, mas desenvolveram algum raciocínio pessoal ou indicaram causas e/ou fatores foram valoradas.
As causas apontadas com maior frequência para a violência foram : a ausência de educação e o desemprego, com soluções correlatas de educação profissional e de qualidade e abertura de postos de trabalho.
A questão das drogas também foi bastante mencionada, e a solução correlata mais indicada foi o acolhimento de usuários de drogas.
Outro tema bastante citado foi a dissolução das famílias e dos valores familiares.
Tendo em vista que o texto falava de moradores de rua, a ausência de moradia apareceu em muitas respostas, bem como a disputa territorial e a necessidade de sobrevivência.
Nesse sentido, as propostas apontavam para a necessidade de ajuda da área de Serviço Social das prefeituras e atendimento psicológico gratuito.
Numa ótica mais repressora, surgiram críticas quanto à leniência das leis, acompanhadas da indicação da necessidade de atuação mais rigorosa do Poder Judiciário, como também da aplicação de penas maiores, negativa de esmolas, proibição do trabalho informal ( guardador de carros ), dentre outras.
Algumas respostas indicaram como causa o modelo capitalista, o consumismo e a ganância.
Outras indicaram o trabalho infantil e a correlata necessidade de aplicação do ECA ( Estatuto da Criança e do Adolescente ).
Em relação ao tema violência, inúmeras respostas apresentaram visões de intolerância quanto aos menores infratores, postulando a diminuição da idade para efeitos de criminalização, bem como defendendo a majoração de penas.
Essa perspectiva revela uma cisão social, de certa forma estimulada pelo tipo de divulgação da mídia quanto aos temas em foco, o que, sem dúvida, ficou evidenciado em algumas respostas.
A existência de inúmeras respostas pugnando pela maior repressão estatal e pelo aumento das penas, acrescido da busca pela diminuição da idade penal, revela que muitos estudantes não foram atingidos por uma formação superior que se exige humanista, ética e comprometida socialmente, que busque os fundamentos para resolver o problema da violência por meio de práticas democráticas e de inclusão social.
As respostas desse tipo também revelaram um alto grau de intolerância para com os moradores de rua, associando -os muitas vezes ao tráfico de drogas, à exploração indevida do espaço público, bem como a crimes como extorsão dos motoristas quando pedem dinheiro para guardar os veículos.
Nessa trilha, verifica -se lacuna na discussão de temas importantes para a formação de profissionais de nível superior tais como : sociodiversidade, multiculturalismo e violência ; tolerância/intolerância ; inclusão/exclusão.
Para aqueles que responderam em conformidade com o padrão de resposta verificou -se uma plena formação social, conclamando pelo auxílio não só do Estado, mas também da Sociedade no amparo daqueles que não têm as mínimas condições de sobrevivência e que estão nas ruas. 3.3.1.5 Análise de Língua Portuguesa das Questões Discursivas do Componente de Formação Geral Os dados de Matemática, obtidos a partir das respostas às questões discursivas do Componente de Formação Geral no que tange à Língua Portuguesa, encontram -se na Tabela 3.16 e no Gráfico 3.10.
Nesse aspecto, os estudantes, de todo o Brasil, tiveram média 52,6.
A maior média com respeito à Língua Portuguesa foi obtida na região Sul ( 54,5 ), e a menor, na região Norte ( 51,0 ).
Quanto à dispersão das notas, o desvio padrão de todo o Brasil foi 24,0.
O menor desvio padrão foi obtido na região Norte ( 22,9 ) e o maior desvio padrão na região Centro-Oeste ( 24,6 ).
A mediana das notas de Língua Portuguesa foi 60,0 para o Brasil como um todo, a mesma obtida nas regiões Nordeste e Centro-Oeste.
Nas demais regiões a mediana foi : 57,5 na região Norte ; e 62,5 nas regiões Sudeste e Sul.
A nota máxima para todo o Brasil foi de 95,0, com, pelo menos, um estudante tirando essa nota nas regiões Norte e Sudeste, enquanto nas demais regiões, a nota máxima foi : 92,5 nas regiões Nordeste e Sul ; e 90,0 na região Centro-Oeste.
Além disso, a nota mínima foi zero em todas as regiões do país, sem exceção.
O Gráfico 3.10 mostra a distribuição das notas de Língua Portuguesa do Componente de Formação Geral.
Observa -se que a maior frequência corresponde aos estudantes que deixaram ambas as questões discursivas em branco, para o Bacharelado, e aos que obtiveram nota no intervalo ( 60 ; 70 ], para a Licenciatura. 3.3.1.6 Comentários sobre a correção das respostas de Formação Geral com respeito à Língua Portuguesa Ao encaminhar as questões 1 e 2 na direção da produção de um texto dissertativo, esperava -se que o estudante utilizasse seus conhecimentos sobre o assunto e estruturasse seus textos de acordo com as exigências do registro formal próprio dessa situação comunicativa.
Essa configuração determina exigências quanto : à adequação da seleção vocabular, ao desenvolvimento do conteúdo, à estruturação sintática dos períodos, à organização lógica das ideias, à utilização de procedimentos de encadeamento textual e referenciação, à obediência às exigências morfossintáticas próprias da modalidade escrita da norma-padrão, ao respeito às regras ortográficas e às regras de acentuação gráfica.
O padrão de respostas utilizado na avaliação considerou os aspectos relevantes ao bom desempenho linguístico como competências distintas, de modo a permitir um mapeamento detalhado do domínio dos recursos disponíveis na Língua Portuguesa para a comunicação escrita formal.
Com base nesse objetivo, foram avaliados os seguintes aspectos : a ) Estruturação textual condizente com o gênero solicitado e o modo de organização textual expositivo adequado ao gênero – essa competência envolve : a estruturação sintática condizente com o padrão da modalidade escrita formal da língua portuguesa, de modo a garantir a clareza necessária ; a distribuição do conteúdo do texto em parágrafos, de modo a garantir a sua organização temática ; a utilização de operadores discursivos que contribuam para a progressão temática do texto, estabelecendo relações lógicas entre as ideias apresentadas, tanto do ponto de vista intrafrasal, como do interfrasal ; a utilização de procedimentos de referenciação lexical e pronominal que permitam a retomada de referentes textuais ; o respeito às regras de pontuação como fator de estruturação do período.
Espera -se, portanto, que o estudante recorra a procedimentos linguístico-discursivos para organizar seu texto, permitindo o encadeamento lógico entre suas partes de forma a garantir a progressão e a coerência textuais.
Isso significa que os seguintes procedimentos foram penalizados, de acordo com o padrão de respostas proposto :  a estrutura lógico-gramatical do texto fica comprometida pela elaboração de frases fragmentadas ;  sequência justaposta de ideias sem encaixamentos sintáticos, reproduzindo hábitos da oralidade ;  elaboração de frase com apenas oração subordinada, sem oração principal ;  emprego equivocado do conector ( preposição, conjunção, pronome relativo, alguns advérbios e locuções adverbiais ) que não estabeleça relação lógica entre dois trechos do texto e prejudique a compreensão da mensagem ;  emprego do pronome relativo sem a preposição, quando obrigatória ;  repetição ou substituição inadequada de palavras sem utilização dos recursos oferecidos pela língua ( pronome, advérbio, artigo, sinônimo ) ;  emprego inadequado dos pronomes relativos “ cujo(a) ” e “ onde ” ;  utilização inadequada dos sinais de pontuação que comprometa a clareza textual. b ) Respeito às convenções ortográficas da norma-padrão da Língua Portuguesa – essa competência envolve o domínio das regras de acentuação gráfica e da grafia padrão das palavras ( com ausência de abreviaturas próprias da linguagem da internet ), de acordo com as convenções estabelecidas pela legislação em vigor e consubstanciadas no Vocabulário Ortográfico da Língua Portuguesa, editado pela Academia Brasileira de Letras ( com aceitação da legislação anterior, no caso das regras relativas ao uso do hífen e da acentuação gráfica ).
Espera -se que o participante :  grafe corretamente as palavras ;  respeite as regras de acentuação gráfica ;  empregue maiúsculas em início de frase, em nomes próprios de pessoas, lugares ou instituições ;  evite abreviações como p/, vc, tb, pq, tá, né, usadas muitas vezes em escrita informal e na internet ;  obedeça às regras de separação de sílabas no final da linha. c ) Domínio dos diferentes aspectos morfossintáticos próprios da modalidade escrita formal da norma-padrão da Língua Portuguesa – essa competência envolve : a concordância nominal, a concordância verbal, a regência nominal, a regência verbal, a flexão nominal, a flexão verbal, a correlação entre os tempos verbais, a colocação pronominal e a utilização de sinais de pontuação que contribuam para a organização lógica da frase e do texto.
Espera -se que o participante :  flexione o verbo para estabelecer concordância de número com o sujeito da frase ;  flexione o artigo, o adjetivo e o pronome para concordar em número e em gênero com o substantivo a que se referem ;  observe a regência nominal e a verbal, utilizando a preposição adequada depois de um substantivo, um verbo ou um adjetivo ;  empregue adequadamente o acento grave indicador de crase nos casos em que se fizer necessário ;  obedeça às regras de colocação pronominal ( próclise e ênclise ), distintas dos hábitos da oralidade ou da escrita informal ; d ) Seleção vocabular adequada à modalidade escrita formal da Língua Portuguesa, exigida pela situação comunicativa – essa competência envolve a precisão na utilização do vocabulário relacionado à temática solicitada pela questão ; a ausência de marcas da oralidade, como termos de sentido de muito genérico ( “ coisa ”, “ negócio ”, “ você ” ) e termos de registros mais informais ( como gírias, jargões, frases feitas, ditados populares, termos regionais ).
Assim, espera -se que o participante respeite a adequação vocabular não utilizando gírias ou expressões coloquiais, evite repetição desnecessária de palavras e utilize um vocabulário mais formal, adequado ao texto de caráter dissertativo.
A escolha dessas competências para subsidiar o processo de avaliação apoia -se na concepção de que, no desempenho dos graduandos, a modalidade escrita tem apresentado uma intensa simplificação, originada no padrão da modalidade oral da Língua Portuguesa.
No caso do texto de base dissertativa, inscrito em um registro formal, a distância entre as duas modalidades é ainda maior, o que provoca situações de hipercorreção ( desvios provocados pela incorporação indevida de uma regra da norma-padrão ) e de truncamentos sintáticos ( estruturas frasais incompreensíveis devido à complexidade sintática própria da modalidade escrita ).
Observam -se, então, os seguintes aspectos que marcam essa distinção entre as duas modalidades, devido à excessiva simplificação da modalidade falada : a ) redução drástica de estruturas subordinadas, compensada pelo aumento na frequência de estruturas coordenadas e absolutas, por um lado, ou pela elaboração de estruturas truncadas pelo excesso de ideias sem a devida conexão subordinativa ; b ) redução no uso de conectores para expressar relações lógicas essenciais à construção do texto, substituídas pela exigência de inferência por parte do interlocutor para suprir a sua ausência ; c ) redução cada vez maior do uso do subjuntivo, ao lado da ampliação do uso do indicativo combinado a estruturas frasais coordenadas ou absolutas ; d ) empobrecimento do processo de referenciação, com a repetição exaustiva de pronomes ou nomes ; e ) simplificação extrema da marcação da categoria tempo na morfologia verbal ; f ) falta de domínio de vocabulário mais abstrato e de maior complexidade, essencial ao desenvolvimento do processo dissertativo ; g ) redução drástica no emprego da acentuação gráfica, processo intensificado pela divulgação imprecisa das mudanças promovidas pelo último acordo ortográfico.
Os aspectos macroestruturais da elaboração do texto não foram avaliados neste processo, para não penalizar duplamente os estudantes, já que a banca de formação geral, composta por profissionais de diferentes áreas do conhecimento, encarregou -se da avaliação do conteúdo desenvolvido nas questões.
São eles : progressão temática, coerência na relação com os conhecimentos de mundo dos usuários da língua, inteligibilidade, atendimento ao solicitado no enunciado do ponto de vista do desenvolvimento do conteúdo, entre outros.
A grade de avaliação do desempenho linguístico considerou, portanto, três grandes grupos de competências, segundo os aspectos explicitados anteriormente : 1.
Aspectos ortográficos : domínio das convenções ortográficas : grafia de vogais e consoantes, uso de maiúsculas e minúsculas, emprego do hífen e acentuação gráfica ; 2.
Aspectos textuais : domínio dos procedimentos de estruturação textual do ponto de vista microestrutural : organização interna dos períodos, emprego de conectores para a articulação lógica entre os períodos e entre os parágrafos, emprego de marcas de referenciação lexical e pronominal ; utilização dos sinais de pontuação que contribuem para a organização lógica da frase. 3.
Aspectos morfossintáticos e vocabulares : domínio das regras de caráter morfossintático estabelecidas como modelares do ponto de vista da modalidade escrita formal da norma-padrão da Língua Portuguesa : concordância nominal e verbal, regência nominal e verbal, colocação pronominal, flexão nominal e verbal, correlação entre tempos e modos verbais, ausência de marcas de oralidade.
A seleção vocabular adequada à modalidade escrita formal da Língua Portuguesa foi incorporada a essa última competência, tendo em vista a intersecção entre as duas do ponto de vista das exigências do registro formal da modalidade escrita da norma-padrão.
Os resultados da avaliação correspondem aos seguintes aspectos observados em cada competência : Aspectos ortográficos : A correção foi realizada classificando os textos em cinco níveis, nível zero a quatro.
O desempenho dos estudantes ficou concentrado, majoritariamente, no nível 3 ( de 4 a 7 desvios ), devido, principalmente, aos desvios de acentuação.
O número de desvios de grafia foi reduzido.
Diferentemente do resultado do ENADE/2013, aumentou a porcentagem de textos sem desvios ( enquadrados no nível 4 ) e diminuiu a porcentagem de textos enquadrados no nível 2.
Alguns textos, com pior desempenho, foram enquadrados no nível 1 ( de 8 a 12 desvios ), enquanto o nível zero foi atribuído a pouquíssimos casos, já que ele revela falta absoluta de domínio das convenções ortográficas.
Observou -se, portanto, que existe uma diferença muito grande de desempenho nos dois aspectos analisados : baixo índice de desvios ortográficos e grande índice de desvios de acentuação.
Em vários casos, ocorre ausência completa de acentuação gráfica.
Os resultados revelam, portanto, que a tendência dominante entre os universitários brasileiros é a eliminação da acentuação gráfica, provavelmente motivada pela vivência dos jovens relacionada aos aplicativos de comunicação via internet ( redes sociais e e-mails ).
Nesse tipo de comunicação, devido ao ritmo intenso de troca de mensagens, o uso de acentos gráficos foi praticamente abolido.
Outro fator que pode ter relação com essa tendência é a ausência de esclarecimento dos meios de comunicação, das autoridades e das escolas sobre as decisões do Acordo Ortográfico da Língua Portuguesa de 1990, gerando um estado de indefinição para os estudantes.
Os casos mais sistemáticos de eliminação do acento indicador da sílaba tônica são : • palavras proparoxítonas ( “ ridiculos ”, “ publicas ”, “ lideres ”, “ politicos ” ) ; • palavras paroxítonas terminadas em ditongo crescente ( “ homicidios ”, “ latrocinio ”, “ individuo ”, “ dependencia ” ) ; • palavras oxítonas ( “ ninguem ”, “ esta ”, “ ate ”, “ ai ” ).
Por outro lado, destaca -se o uso indevido do acento gráfico em determinadas palavras, como observado nas grafias * “ jornáis ”, * “ telejornaís ”, * “ propíciar ”, * “ medídas ”, * “ dígnidade ”, * “ cídades ”.
Quanto ao domínio das convenções relativas à grafia das palavras, observam -se desvios como : a hipercorreção pela escolha de “ e ” no lugar de “ i ”, por influência de hábitos da oralidade ( * “ enumeros ” por “ inúmeros ”, * “ entevenção ” por “ intervenção ” ) ; a eliminação do “ r ” marcador do infinitivo verbal ( * “ esta ” no lugar de “ estar ” ).
Outros casos de desvios de grafia relacionados à variação diastrática podem ser observados em *“estrupo ”, *“automovís ”, *“viensse ”, *“camihada ”.
São muito frequentes os seguintes desvios de caráter ortográfico, com repercussão morfossintática : • eliminação da marca de infinitivo ( -r- ) e substituição por acento agudo ( “ está ” no lugar de “ estar ” ) ; • confusão entre “ ão ” e “ am ” nas formas verbais ( “ invadão ” no lugar de “ invadam ” e “ estam ” no lugar de “ estão ” no presente do indicativo ; “ estaram ” no lugar de “ estarão ” no futuro do indicativo ) ; • confusão entre a grafia do verbo “ haver ” ( “ há ” ) e o artigo definido ou a preposição “ a ” ; • uso de hífen para separar pronome átono – tanto uso indevido quanto omissão ( no pretérito imperfeito e futuro do subjuntivo : “ evitar -mos ” no lugar de “ evitarmos ”, “ percebesse ” no lugar de “ percebe -se ” e vice-versa ) ; Observam -se, também, muitos casos de inadequação no uso da maiúscula : ausência de diferença entre a primeira letra e as outras, em início de período, principalmente dos estudantes que adotam a escrita em letra de imprensa ; utilização de maiúscula para destacar determinadas palavras-chave do texto, como “ Violência ”, “ Brasileiros ”, “ Fatores ”, “ Ozônio ”, “ Sustentável ”.
Vale observar, também, que, ao contrário do que se esperava, não apareceram abreviaturas próprias do “ internetês ”, ou seja, dos hábitos de comunicação escrita adquiridos pelo uso de redes sociais e emails.
Aspectos textuais : Quanto a esses aspectos a correção também classificou os textos em cinco níveis ( zero a quatro ), em função da quantidade de erros apresentados.
O desempenho dos estudantes ficou concentrado, majoritariamente, nos níveis 3 e 2, devido à grande ocorrência de problemas de estruturação textual.
Foi muito baixo o número de textos que não apresentaram qualquer problema estrutural e, portanto, ficaram enquadrados no nível 4.
Alguns textos com pior desempenho foram enquadrados no nível 1, enquanto o nível zero foi atribuído a pouquíssimos casos, já que ele revela a existência de um texto sem articulação e com comprometimento do sentido.
Observou -se que a grande maioria dos estudantes não distribuiu as ideias em parágrafos, talvez devido ao pequeno número de linhas disponibilizadas para a resposta da questão ou, quem sabe, pela suposição de que não seria necessária essa divisão por não se tratar de um texto no modelo de uma redação dissertativo-argumentativa, como solicitado nos vestibulares.
Outro aspecto observado na estruturação textual foi a divisão em dois itens, provavelmente motivados pelo encaminhamento do enunciado das questões, que apresentavam dois pontos a serem detalhados : a questão 1 solicitava que os estudantes dissertassem sobre as consequências do transporte motorizado ( a ) e apresentassem ações de intervenção por parte do poder público ( b ) ; a questão 2 solicitava que os estudantes dissertassem sobre as causas da violência urbana ( a ) e os fatores para evitá -la ( b ).
Essa competência pode ser considerada como a mais problemática no que diz respeito ao desempenho linguístico dos estudantes, porque são muitos os problemas observados, desvios acumulados durante toda a formação escolar e que não se resolvem com um estudo autodidata, como acontece com regras ortográficas ou morfossintáticas : sequência justaposta de ideias sem encaixamentos sintáticos ; redução drástica de estruturas subordinadas, ao lado do aumento na frequência de estruturas coordenadas e absolutas ; redução no uso de conectores para expressar relações lógicas essenciais à construção do texto, substituídas pela exigência de inferência por parte do interlocutor para suprir a sua ausência ; emprego equivocado de operadores que não estabelecem relações lógicas coerentes entre ideias do texto ; emprego inadequado do pronome relativo ( com omissão da preposição ou a utilização de pronome inadequado, como “ onde ” ) ; repetição exaustiva de termos sem a utilização de procedimentos mais sofisticados de substituição ( hiperonímias, hiponímias, nominalizações, expressões metafóricas ) ; frases fragmentadas que comprometem a estrutura lógico-gramatical ; frases formadas apenas por oração subordinada, sem oração principal.
Um importante aspecto a destacar é o baixíssimo desempenho de uma parte dos estudantes em relação à estrutura formal do texto produzido, o que é extremamente preocupante ao se levar em conta que são graduandos em fase final de formação.
São frequentes os casos de desvios de estruturação frasal, com uso inadequado ou ausência de conectivos entre parágrafos e entre frases.
Em uma parte dos textos, falta um mínimo de textualidade e de domínio do registro padrão da língua.
Na verdade, observam -se relações linguísticas quase agramaticais, como as estabelecidas pela sequência de gerúndios sem o apoio de um ponto de partida para a organização das informações gramaticais e semânticas, ou seja, sem uma oração principal.
Quanto à utilização dos mecanismos de referenciação, deve -se destacar a ocorrência, em uma boa parte dos textos, de repetições de palavras ou expressões sem a utilização de termos sinônimos ou pronomes, como seria adequado.
Quanto à utilização dos sinais de pontuação, observou -se uma grande precariedade nos textos analisados.
É muito frequente a ocorrência de parágrafos sem marca interna de pontuação para separar os períodos.
Vale observar que não foi penalizada a ausência de vírgula para destacar locuções ou adjuntos adverbiais de pequena extensão deslocados de posição na frase, por ser um uso opcional.
São os seguintes os tipos de problemas encontrados :  vírgula : utilização de vírgula para separar o sujeito e o predicado ; ocorrência de apenas uma das vírgulas para separar uma palavra, uma expressão ou uma oração encaixada ; uso de vírgula no lugar do ponto para separar ideias que constituem períodos distintos ; ausência de vírgula para separar elementos de uma enumeração ; ausência de vírgula para separar oração adjetiva explicativa ou utilização inadequada para separar oração adjetiva restritiva ;  ponto e vírgula : utilização do ponto e vírgula no lugar de vírgula ;  ponto final : ausência de ponto final para separar períodos.
Aspectos morfossintáticos e vocabulares : Da mesma forma que nos aspectos anteriores os textos foram classificados em níveis de zero a quatro.
O desempenho dos estudantes ficou concentrado, majoritariamente, nos níveis 3 e 2, devido à grande ocorrência de problemas de regência e concordância.
O nível 4 foi atribuído a um número menor de textos.
Alguns textos com pior desempenho foram enquadrados no nível 1, enquanto o nível zero foi atribuído a pouquíssimos casos, já que ele revela a existência de um texto sem o respeito às mínimas exigências morfossintáticas da norma-padrão e com comprometimento do sentido.
Os resultados são muito transparentes em relação aos aspectos mais problemáticos do desempenho dos estudantes.
O desvio mais frequente, em relação à regência, é a falta do sinal indicativo da crase – isso revela que o usuário não tem consciência de que, sob a forma do termo “ a ”, existe a presença de uma preposição “ a ”, exigida pela regência do termo anterior.
Embora em outros exames, como o Enem, a falta de crase seja penalizada em acentuação, nesta correção esse desvio foi considerado no âmbito dos aspectos morfossintáticos.
Outro problema relacionado à regência verbal e à nominal, encontrado frequentemente nas questões, foi a ausência de preposição antes de pronome relativo, processo generalizado na modalidade oral da língua, em situações de registro informal.
Apesar da possibilidade de que essa alteração de regência se generalize no padrão escrito da Língua Portuguesa, como já está ocorrendo até em textos jornalísticos, o não emprego da preposição foi penalizado neste processo de avaliação.
Outro desvio muito frequente diz respeito aos processos de concordância verbal e de concordância nominal.
Quanto à concordância de número, observou -se ausência de marca ( com sujeito anteposto ou posposto ) ou uso indevido ( uso inadequado da marca de plural comandado pelo núcleo plural da locução adjetiva, apesar de o substantivo que funciona como núcleo do sintagma nominal estar no singular ).
Uma ocorrência que se destacou foi a ausência de acento circunflexo na forma plural do presente do indicativo dos verbos “ ter ” e “ vir ”, que foi considerada como um desvio na concordância verbal e não na acentuação gráfica.
Quanto à concordância de gênero, vários casos foram observados, normalmente no âmbito de sintagmas nominais longos, em que o adjetivo está afastado do substantivo. concordância de gênero, vários casos foram observados, normalmente no âmbito de sintagmas nominais longos, em que o adjetivo está afastado do substantivo.
Deve -se destacar uma ocorrência não observada no ENADE/2013 : o aparecimento da marca de plural em verbos ou adjetivos relacionados a núcleos substantivos no singular, evidenciando um processo de hipercorreção.
Quanto à questão da colocação pronominal, foram poucos os casos observados.
Apesar de serem aspectos relacionados à oralidade, concluiu -se que, no registro escrito formal, a maioria dos estudantes já incorporou regras como a não introdução da frase por um pronome oblíquo e a próclise na presença de um termo atrator.
Não se adotou, entretanto, o padrão excessivamente formal descrito pelas gramáticas normativas em relação à posição do pronome oblíquo em locuções verbais, já que esse uso está muito distante da prática cotidiana, até em textos mais formais.
Quanto aos aspectos vocabulares, alguns tipos de inadequação foram observados : expressões da oralidade apareceram em algumas respostas, mas sem maior relevância do ponto de vista quantitativo ; seleção vocabular incompatível com o contexto, gerando situações de falta de inteligibilidade ; falta de domínio de vocabulário mais abstrato e de maior complexidade, essencial ao desenvolvimento do texto de base dissertativa.
Várias marcas de oralidade foram identificadas, embora não em alta frequência : o uso do pronome relativo “ onde ” como relativo universal, falta de artigo definido antes de substantivo, repetição de palavras por falta de vocabulário, reduções como “ tá ”, “ pra ”, “ pro ”, “ prum ”, expressões informais. 3.3.2 Componente de Conhecimento Específico Na parte da prova relativa às questões discursivas no Componente de Conhecimento Específico ( Tabela 3.17 ), observa -se que a média foi bem mais baixa do que para as questões discursivas do Componente de Formação Geral.
Enquanto no Componente de Formação Geral a média para estudantes de Matemática ( Bacharelado e Licenciatura ) de todo o Brasil foi 49,3, na parte de Conhecimento Específico a média foi 10,7.
A maior média deste componente foi obtida pelos estudantes da região Sudeste ( 13,2 ), e a menor, pelos da região Norte ( 6,3 ).
Quanto à variabilidade das notas, o desvio padrão de todo o Brasil foi 16,6.
O maior desvio padrão foi encontrado na região Sudeste ( 19,4 ), e o menor, na região Norte ( 10,8 ).
As notas máximas ( 100,0 ) e as notas mínimas ( 0,0 ) foram as mesmas em todas as regiões do Brasil, sem exceção.
A mediana do Brasil como um todo foi 5,0, a mesma obtida nas regiões Sudeste e Sul.
Nas demais regiões a mediana foi : 0,0 na região Norte ; 3,3 na região Nordeste ; e 5,9 na região Centro-Oeste.
O Gráfico 3.11 representa a distribuição das notas nas questões discursivas no Componente de Conhecimento Específico segundo a opção de Bacharelado/Licenciatura.
As modas destas distribuições ocorrem no primeiro intervalo, [ 0 ; 10 ].
O percentual de estudantes de Licenciatura com notas no intervalo [ 0 ; 10 ] é muito maior do que o de Bacharelado, enquanto esta proporção se inverte a partir do intervalo ( 20 ; 30 ].
A análise de cada uma destas questões será feita a seguir. 3.3.2.1 Análise da Questão Discursiva 3 do Componente de Conhecimento Específico A questão 3 foi respondida por concluintes do Bacharelado e da Licenciatura e os resultados aferidos encontram -se descritos na Tabela 3.18.
O desempenho dos estudantes de todo o Brasil nesta questão foi o mais baixo dentre as três questões discursivas desse componente, e a média dos estudantes de todo o Brasil foi 5,9.
A menor média nessa questão foi obtida pelos estudantes da região Norte ( 2,7 ), enquanto a maior média foi obtida na região Centro-Oeste ( 8,3 ).
Quanto à variabilidade das notas, o desvio padrão de todo o Brasil foi 16,6.
O maior desvio padrão foi obtido na região Sudeste ( 19,4 ), enquanto o menor foi obtido na região Norte ( 10,1 ).
A nota máxima, 100,0 pontos, foi alcançada por, pelo menos, um estudante em todas as regiões do Brasil, sem exceção.
As medianas e as notas mínimas foram zero em todas as regiões do país, o que significa que pelo menos 50 \% dos estudantes de todo o Brasil obtiveram nota zero ou deixaram a questão em branco.
O Gráfico 3.12 mostra a distribuição das notas na questão discursiva 3 do Componente de Conhecimento Específico da área de Matemática para estudantes do Bacharelado e da Licenciatura.
Essa distribuição tem moda nas respostas em branco, com percentual muito maior para os estudantes de Licenciatura. 3.3.2.2 Comentários sobre as respostas à Questão Discursiva 3 A questão abordava tópicos de Álgebra Linear e de Geometria Analítica, todos bastante relevantes na formação do bacharel e do licenciado em Matemática.
O conteúdo abordado na questão costuma ser apresentado em disciplinas básicas de Geometria Analítica e Álgebra Linear, oferecidas na primeira metade dos cursos de Bacharelado e de Licenciatura em Matemática.
A questão foi minimamente contextualizada, e o que foi solicitado aos estudantes nos três itens articulou -se pouco com o contexto apresentado, sobretudo no que se refere ao item ‘ c ’.
De modo geral, o espaço destinado à resolução da questão foi adequado, mas, dependendo das soluções dadas pelos estudantes, pode ter sido justo, ou insuficiente.
A questão pode ser considerada de dificuldade média.
O item ‘ a ’ exigiu do estudante a compreensão de uma propriedade geométrica da elipse e/ou o encaminhamento algébrico necessário à obtenção de sua equação geral, a partir do conhecimento da medida de seu eixo maior e das coordenadas dos seus focos.
Foi o item que mais exigiu do estudante articulações de conteúdos de geometria analítica. \textbf{Grosso} modo, o padrão de respostas adotado considerou que o item ’ a ’ poderia ser resolvido de duas maneiras :  confirmando a propriedade fundamental da elipse acerca das distâncias do ponto p(0,2 + √3 ) aos seus focos ( F₁ e F₂ ), a saber, d(p, F₁ ) + d(p, F₂ ) = 2a, onde 2a indica o comprimento do seu eixo maior, ou  determinando a medida do semieixo menor ( b ) para, a partir da relação a² = b² + c², obter a equação da elipse e verificar que o ponto p(0,2 + √3 ) a satisfazia.
O item ‘ b ’ apresentou um operador linear específico e solicitou a obtenção das imagens dos focos da elipse F₁ ( -1,2 ) e F₂ ( 1,2 ).
O item ‘ b ’ foi, certamente, o mais elementar dentre os três itens apresentados na questão, uma vez que, para respondê -lo, os estudantes deveriam apenas compreender o conceito de função.
Nesse ponto, o padrão de respostas considerou o único encaminhamento esperado : a substituição das coordenadas dos focos na expressão algébrica do operador e a seguinte obtenção das imagens.
Cabe considerar que, para efeitos \textbf{avaliativos}, o item ’ b ’ foi redundante ao solicitar que fizessem a mesma coisa duas vezes, uma vez que o que difere as obtenções de T(1,2) e de T(-1,2) são detalhes periféricos diante do que se desejava avaliar.
A obtenção de apenas uma imagem já seria suficiente para os estudantes manifestarem seu conhecimento nesse sentido.
Por fim, o item ‘ c ’ solicitou os autovalores do operador linear dado, tendo sido considerados, no padrão de respostas, dois possíveis encaminhamentos distintos, ainda que matematicamente equivalentes.
O primeiro consistiu na imediata obtenção das raízes do polinômio característico p(λ) = det(A − λI ) e, o segundo, na utilização da definição de autovalor e autovetor, T(v) = λ·v, a fim de se obter um sistema de equações a ser resolvido em λ.
De modo geral, as respostas corretas dadas pelos estudantes se detiveram nos casos previstos no padrão de respostas.
No item ‘ a ’, alguns estudantes encaminharam uma solução que, de certa forma, colocou -se entre as duas possibilidades presentes no padrão.
Após perceberem que o eixo maior da elipse era horizontal, que o centro da elipse era o ponto ( 0,2 ) e que b = √3, concluíram que o ponto ( 0,2 + √3 ) deveria pertencer à elipse.
Um erro comum ocorreu pela consideração de que a igualdade d(p, F₁ ) = d(p, F₂ ) era suficiente para garantir que o ponto p pertencia à elipse.
No problema apresentado, de fato, tal igualdade teria sido suficiente, caso os estudantes estivessem utilizado o fato de que ambas eram iguais a 2, a fim de concluir que d(p, F₁ ) + d(p, F₂ ) = 4 = 2a.
Infelizmente, a maior parte julgou que a simples igualdade era suficiente para a conclusão do item ‘ a ’, sendo desconsiderada a relevância de que ambas as distâncias eram iguais ao comprimento do semieixo maior da elipse.
A solução do item ‘ b ’ seguiu o padrão de respostas esperado, dado pela substituição das coordenadas dos focos na expressão algébrica do operador T.
A resolução correta do item ‘ c ’, também, se deu conforme o previsto no padrão de respostas.
Esse foi o item no qual os estudantes mais erraram, ou sequer esboçaram uma solução.
Notando que o operador T proporcionava o caso mais simples possível para a determinação de autovalores reais, fica corroborada a hipótese de que os estudantes não conheciam o conceito de autovalor.
Os estudantes oriundos de cursos de Bacharelado tiveram um desempenho significativamente melhor do que aqueles provenientes de cursos de Licenciatura.
A média das notas obtidas pelos estudantes do Bacharelado foi superior à média obtida pelos estudantes de Licenciatura.
A fim de justificar tal diferença de desempenho entre os estudantes dos dois cursos, uma hipótese deve ser considerada.
O conteúdo de Geometria Analítica e Álgebra Linear tratado na questão tem presença consolidada nos cursos de Bacharelado, quase sempre, já nos seus dois períodos iniciais.
Com as A obtenção de apenas uma imagem já seria suficiente para os estudantes manifestarem seu conhecimento nesse sentido.
Por fim, o item ‘ c ’ solicitou os autovalores do operador linear dado, tendo sido considerados, no padrão de respostas, dois possíveis encaminhamentos distintos, ainda que matematicamente equivalentes.
O primeiro consistiu na imediata obtenção das raízes do polinômio característico p(λ) = det(A − λI ) e, o segundo, na utilização da definição de autovalor e autovetor, T(v) = λ · v, a fim de se obter um sistema de equações a ser resolvido em λ.
De modo geral, as respostas corretas dadas pelos estudantes se detiveram nos casos previstos no padrão de respostas.
No item ‘ a ’, alguns estudantes encaminharam uma solução que, de certa forma, colocou -se entre as duas possibilidades presentes no padrão.
Após perceberem que o eixo maior da elipse era horizontal, que o centro da elipse era o ponto ( 0,2 ) e que 𝑏 = √3, concluíram que o ponto ( 0,2 + √3 ) deveria pertencer à elipse.
Um erro comum ocorreu pela consideração de que a igualdade 𝑑(𝑝, 𝐹1 ) = 𝑑(𝑝, 𝐹2 ) era suficiente para garantir que o ponto p pertencia à elipse.
No problema apresentado, de fato, tal igualdade teria sido suficiente, caso os estudantes estivessem utilizado o fato de que ambas eram iguais a 2, a fim de concluir que 𝑑(𝑝, 𝐹1 ) + 𝑑(𝑝, 𝐹2 ) = 4 = 2𝑎.
Infelizmente, a maior parte julgou que a simples igualdade era suficiente para a conclusão do item ‘ a ’, sendo desconsiderada a relevância de que ambas as distâncias eram iguais ao comprimento do semieixo maior da elipse.
A solução do item ‘ b ’ seguiu o padrão de respostas esperado, dado pela substituição das coordenadas dos focos na expressão algébrica do operador T.
A resolução correta do item ‘ c ’, também, se deu conforme o previsto no padrão de respostas.
Esse foi o item no qual os estudantes mais erraram, ou sequer esboçaram uma solução.
Notando que o operador T proporcionava o caso mais simples possível para a determinação de autovalores reais, fica corroborada a hipótese de que os estudantes não conheciam o conceito de autovalor.
Os estudantes oriundos de cursos de Bacharelado tiveram um desempenho significativamente melhor do que aqueles provenientes de cursos de Licenciatura.
A média das notas obtidas pelos estudantes do Bacharelado foi superior à média obtida pelos estudantes de Licenciatura.
A fim de justificar tal diferença de desempenho entre os estudantes dos dois cursos, uma hipótese deve ser considerada.
O conteúdo de Geometria Analítica e Álgebra Linear tratado na questão tem presença consolidada nos cursos de Bacharelado, quase sempre, já nos seus dois períodos iniciais.
Com as Com as alterações curriculares que se deram nos cursos de Licenciatura, a partir da Resolução CNE/CP nº 2, de 19 de fevereiro de 2002, muitas instituições reorganizaram suas grades curriculares e moveram as disciplinas de Geometria Analítica e Álgebra Linear para o segundo ano do curso, mais especificamente no terceiro e/ou quarto períodos e com uma abordagem menos aprofundada.
De modo geral, as respostas que foram consideradas fracas ( notas de zero a 30 ) foram aquelas em que apenas um dos itens foi respondido, quase sempre o item ‘ b ’.
As respostas que foram consideradas medianas ( notas de 35 a 70 ), foram aquelas em que o estudante acertou um dos itens ( quase sempre o item ‘ b ’ ) e acertou parcialmente um segundo ( quase sempre o item ‘ a ’ ).
Algumas respostas medianas consistiram apenas na resolução correta do item ‘ c ’.
As respostas consideradas boas ( notas de 75 a 100 ) foram aquelas em que pelo menos dois itens foram feitos corretamente ( um deles - quase sempre - o item ‘ b ’ e acertou o item ‘ c ’, pelo menos parcialmente ).
Como dito anteriormente acerca do item ‘ a ’, as respostas corretas seguiram, de modo geral, o padrão de respostas determinado.
Um erro cometido com frequência pelos estudantes se deu pela consideração de que a igualdade 𝑑(𝑝, 𝐹1 ) = 𝑑(𝑝, 𝐹2 ) é suficiente para garantir que o ponto p(0,2 + √3 ) pertence à elipse.
No problema apresentado, tal igualdade teria sido, de fato, suficiente, caso tivessem percebido que ambas as distâncias eram iguais a 2 e, por isso, ter -se-ia verificada a igualdade 𝑑(𝑝, 𝐹1 ) + 𝑑(𝑝, 𝐹2 ) = 4 = 2𝑎.
Infelizmente, a maior parte julgou que a simples igualdade era suficiente para a conclusão do item ’ a ’, sendo desconsiderada a relevância de que ambas as distâncias eram iguais ao comprimento do semieixo maior da elipse.
Houve casos em que os estudantes se equivocaram em cálculos/interpretações e obtiveram a equação de outra elipse que, por vezes, tinha p como um dos seus pontos.
Acerca do item ‘ b ’, verificou -se que a maioria dos estudantes que tiveram algum acerto na questão 3, tiveram -no em tal item, pelo menos parcialmente.
Os eventuais erros cometidos no item ‘ b ’ se deram após a substituição das coordenadas dos focos na expressão de T, mais especificamente no trato com os sinais e/ou na realização do cálculo mental.
Muitos obtiveram corretamente a imagem T(1,2), mas se equivocaram nos sinais durante a obtenção da imagem T(-1,2 ).
A resolução correta do item ‘ c ’ se deu conforme o previsto no padrão de respostas.
Os erros cometidos no item ‘ c ’ perpassaram, quase sempre, erros de conta inerentes à determinação do polinômio característico, seja por meio da fórmula envolvendo o determinante, ou do sistema linear obtido a partir da definição de autovetor/autovalor.
Vários estudantes obtiveram polinômios característicos errados, mas calcularam suas raízes e afirmaram que as mesmas eram os autovalores solicitados.
Como dito anteriormente, houve casos em que os estudantes descobriram os autovalores testando valores inteiros para .
Essa é uma prática indesejável e insuficiente para resolver problemas gerais similares.
Após a correção das questões relativas aos estudantes de Licenciatura, verificou -se que acima 65 \% foram deixadas em branco.
Esse alto percentual pode levantar suspeitas acerca do comprometimento dos estudantes na realização da prova do ENADE/2014.
No entanto, como há outras questões em que o percentual de respostas em branco foi significativamente menor ( como na questão 5 – acima de 35 \% ), aparentemente, os estudantes se debruçaram sobre as questões apresentadas.
Como dito anteriormente, a média das notas obtidas pelos de Licenciatura nessa questão foi baixa e menor do que o valor do item ‘ b ’ ( o mais elementar ).
Esse é um dado alarmante, que pode corroborar a hipótese de que os conteúdos abordados pela questão não eram conhecidos pelos estudantes.
Isso pode ter se dado em consequência da não apresentação de tais conteúdos nos cursos e/ou da baixa qualidade dos cursos de Licenciatura, em termos mais amplos.
Essencialmente, os estudantes que conhecessem minimamente o conceito de função e a elipse teriam alcançado um desempenho mediano na questão.
Inúmeras soluções apresentadas mostraram que muitos estudantes sequer souberam como obter a imagem de um ponto do domínio de uma função. 3.3.2.3 Análise da Questão Discursiva 4 do Componente de Conhecimento Específico ( Bacharelado ) A questão discursiva 4 de Conhecimento Específico proposta para os concluintes do Bacharelado foi anulada pela Comissão Assessora.
Os comentários da banca de correção a respeito dessa questão encontram -se na próxima seção. 3.3.2.4 Comentários sobre as respostas à Questão Discursiva 4 do Componente de Conhecimento Específico ( Bacharelado ) Essa questão apresentava um problema elementar acerca da topologia de um espaço métrico específico.
O primeiro item da questão solicitava a demonstração de que tal conjunto é efetivamente um espaço métrico, quando munido de uma função distância fornecida.
Em seguida, em seu segundo item, solicitava a demonstração de que tal espaço métrico é discreto.
Infelizmente, o enunciado da questão apresentava algumas fragilidades.
No conjunto P das “ palavras com letras ” ( conforme referido no item ‘ a ’ ), a distância de Hamming não é uma métrica : o domínio da função fornecida não é P x P, pois d(x,y) não está definida quando as palavras x e y são formadas por um número distinto de letras. “ Casa ” e “ Cadeira ”, por exemplo, são elementos do conjunto P das “ palavras com letras ”, mas a definição fornecida não permite a obtenção de d(Casa,Cadeira ).
O enunciado do item ‘ a ’ é, portanto, falho ao definir o conjunto das palavras consideradas, pois deveria fixar um número natural arbitrário n de letras, a fim de solicitar a demonstração de que a distância de Hamming define uma métrica no conjunto das palavras formadas por n letras..Devido à gravidade de tal falha, a questão foi anulada. 3.3.2.5 Análise da Questão Discursiva 4 do Componente de Conhecimento Específico ( Licenciatura ) A Tabela 3.19 contém as informações relativas à questão 4 do Componente de Conhecimento Específico ( Licenciatura ).
A média geral do Brasil foi 10,2, sendo a menor média registrada na região Norte ( 4,9 ) e a maior, na região Sudeste ( 14,3 ).
A nota máxima, 100,0 pontos, foi alcançada por, pelo menos, um estudante em todas as regiões do Brasil, sem exceção.
As medianas e as notas mínimas foram 0,0 em todas as regiões do país, o que significa que pelo menos 50 \% dos estudantes de todo o país obtiveram nota zero ou deixaram a questão em branco.
O Gráfico 3.13 representa a distribuição de notas na questão discursiva 4, no Componente de Conhecimento Específico ( Licenciatura ).
A maior frequência ocorre na categoria “ em branco ”. 86 Algumas fragilidades.
No conjunto P das “ palavras com letras ” ( conforme referido no item ‘ a ’ ), a distância de Hamming não é uma métrica : o domínio da função fornecida não é P x P, pois d(x,y) não está definida quando as palavras x e y são formadas por um número distinto de letras. “ Casa ” e “ Cadeira ”, por exemplo, são elementos do conjunto P das “ palavras com letras ”, mas a definição fornecida não permite a obtenção de d(Casa,Cadeira ).
O enunciado do item ‘ a ’ é, portanto, falho ao definir o conjunto das palavras consideradas, pois deveria fixar um número natural arbitrário n de letras, a fim de solicitar a demonstração de que a distância de Hamming define uma métrica no conjunto das palavras formadas por n letras.
Devido à gravidade de tal falha, a questão foi anulada. 3.3.2.5 Análise da Questão Discursiva 4 do Componente de Conhecimento Específico ( Licenciatura ) A Tabela 3.19 contém as informações relativas à questão 4 do Componente de Conhecimento Específico ( Licenciatura ).
A média geral do Brasil foi 10,2, sendo a menor média registrada na região Norte ( 4,9 ) e a maior, na região Sudeste ( 14,3 ).
A nota máxima, 100,0 pontos, foi alcançada por, pelo menos, um estudante em todas as regiões do Brasil, sem exceção.
As medianas e as notas mínimas foram 0,0 em todas as regiões do país, o que significa que pelo menos 50 \% dos estudantes de todo o país obtiveram nota zero ou deixaram a questão em branco.
O Gráfico 3.13 representa a distribuição de notas na questão discursiva 4, no Componente de Conhecimento Específico ( Licenciatura ).
A maior frequência ocorre na categoria “ em branco ”. 3.3.2.6 Comentários sobre as respostas à Questão Discursiva 4 ( Licenciatura ) A questão apresentada para os estudantes de cursos de Licenciatura fornecia, textualmente, as instruções para uma construção geométrica simples : a divisão de um segmento dado ( AB ) em partes ( AD e DB ) que definem a proporção áurea.
Os três itens situam os passos para a obtenção da razão áurea AB/AD, tomando por base a construção acima.
A resolução da questão possui, portanto, duas etapas bem definidas : a compreensão das instruções/realização da construção e os encaminhamentos de cálculo envolvendo as medidas dos segmentos determinados, que incluem o uso do Teorema de Pitágoras.
O padrão de respostas, seguido pelos estudantes, baseou -se na elaboração da construção, seguida dos cálculos.
Alguns estudantes não fizeram a construção no espaço destinado à questão, e o utilizaram apenas para a realização dos cálculos.
Alguns podem ter feito a construção no caderno de prova e, a partir dela, apresentado os cálculos.
O padrão de respostas acordado com a Comissão Assessora considerou, de toda forma, as duas etapas acima expostas.
O conteúdo abordado pela questão é, essencialmente, consolidado na Educação Básica, sendo revisitado, por vezes, em disciplinas de Geometria Básica nos cursos de Licenciatura.
De toda forma, os estudantes não costumam receber os passos de uma construção geométrica, sem o apoio de figuras.
Essa falta de hábito foi revelada pela apresentação de inúmeras construções erradas.
Os cálculos, por sua vez, envolviam radicais e ofereceram dificuldades aos estudantes.
Diante disso, pode -se dizer que a questão foi de dificuldade média.
O espaço destinado à resolução da questão foi justo, mas exigiu do estudante alguma organização.
Muitos estudantes apenas realizaram a construção geométrica no espaço destinado à resolução da questão.
Alguns escreveram corretamente as medidas dos segmentos, em função de x, outros nem isso.
Outros apenas realizaram os cálculos no espaço destinado à resolução da questão, sem apresentar a construção geométrica.
Conforme já comentado, eles podem tê -la feito em outro local, a fim de pouparem espaço.
A maioria das resoluções, no entanto, continha a construção acompanhada pelos cálculos, conforme esperado pelo padrão.
Houve casos em que construções equivocadas originaram cálculos indevidos, mas que, ainda que corretamente executados, foram incapazes de alcançar o resultado solicitado.
Houve, também, resoluções em que as construções foram feitas corretamente, mas erros nos cálculos impediram a obtenção da resposta correta.
Portanto, as respostas erradas dos estudantes se deram, tanto na dimensão geométrica ( que envolve a compreensão das instruções apresentadas textualmente ), como na dimensão algébrica ( que depende apenas de manipulações algébricas elementares ).
Nessa questão, o desempenho dos estudantes dos cursos de Licenciatura foi péssimo.
A média ficou em torno de 10,2 pontos, e o percentual de resoluções em branco ficou em torno de 65 \%, em um quadro semelhante àquele apresentado para a questão 3.
Os erros encontrados revelaram com clareza a dificuldade de extraírem informações básicas de textos simples, um improvável desconhecimento do Teorema de Pitágoras e a baixa proficiência em manipulações algébricas elementares envolvendo frações e radicais.
De modo geral, as respostas consideradas fracas ( notas de zero a 30 ) foram aquelas que trouxeram construções geométricas incompletas, por vezes nem mesmo acompanhadas de medidas.
As respostas consideradas medianas ( notas de 35 a 70 ) foram aquelas que trouxeram construções e/ou cálculos corretos, em um ou dois itens ( normalmente nos itens ’ a ’ e ‘ b ’ ).
As respostas consideradas boas ( notas de 75 a 100 ) foram aquelas em que a figura foi feita corretamente, assim como os cálculos seguintes, exceto, eventualmente, por pequenos erros nas manipulações algébricas.
Na dimensão geométrica, os erros frequentes foram decorrentes da incompreensão das instruções fornecidas e das consequentes falhas na elaboração da construção, assim como na aplicação do Teorema de Pitágoras.
Muitos estudantes fizeram a construção e explicitaram corretamente as medidas envolvidas, mas foram incapazes de aplicar o Teorema de Pitágoras corretamente ( item ’ a ’ ), a fim de obterem outras medidas importantes.
Na dimensão algébrica, os erros foram decorrentes da pouca proficiência no trato com radicais.
Muitos consideraram propriedades falsas como √𝑏2 + 𝑐2 = 𝑏+𝑐, ou cometeram erros básicos nas simplificações/racionalizações.
Um erro bastante comum foi aquele em que o estudante não desenhou o triângulo retângulo descrito e afirmou que AE = EC = 𝑥/2, assumindo o ponto E como o ponto médio do segmento AC.
Com isso, obteve AC = AB = 𝑥 e AD = AE = 𝑥/2 e, consequentemente, concluiu que o número de ouro é igual a AB/AD = 2.
Infelizmente, as respostas obtidas indicam que, de modo geral, os estudantes tiveram não apenas dificuldade em compreender as instruções de encaminhamento da construção, mas também na aplicação do Teorema de Pitágoras ( conhecimento que já deveria estar consolidado desde o Ensino Médio ) e na realização de manipulações algébricas elementares envolvendo radicais.
Ao contrário do que aconteceu na questão 3, não há dúvida de que os conteúdos abordados deveriam ser minimamente conhecidos pelos estudantes ao ingressarem no curso superior.
Tal constatação não dá margem a dúvidas : os estudantes estão ingressando no Ensino Superior sem os conhecimentos mínimos que deveriam ter sido consolidados na Educação Básica e não encontram, nos cursos de graduação em Matemática, políticas de acolhimento capazes de ajudá -los a superarem tais dificuldades. 3.3.2.7 Análise da Questão Discursiva 5 do Componente de Conhecimento Específico ( Bacharelado ) A Tabela 3.20 contém as informações relativas à questão 5 do conjunto do Componente de Conhecimento Específico ( Bacharelado ).
O desempenho dos estudantes nessa questão foi superior ao da questão 3.
A nota média dos estudantes de todo o Brasil foi 43,2.
A maior média foi registrada na região Sudeste ( 51,4 ), enquanto a menor média foi registrada na região Sul ( 31,5 ).
Quanto à variabilidade das notas, o desvio padrão dos estudantes do Brasil, como um todo, foi 36,3.
Enquanto o maior desvio foi encontrado na região Norte ( 43,6 ), o menor foi encontrado na região Centro-Oeste ( 31,4 ).
A nota máxima ( 100,0 ) e a nota mínima ( 0,0 ) são as mesmas em todas as regiões, sem exceção.
A mediana foi 30,0, o que também ocorreu para as regiões Nordeste e Sul.
Para as demais regiões do Brasil a mediana foi 40,0.
O Gráfico 3.14 representa a distribuição das notas da questão discursiva 5 no Componente de Conhecimento Específico ( Bacharelado ).
Mais uma vez, destaca -se o grande número de estudantes que deixaram a questão em branco, quase 1/4.
A moda ocorreu no intervalo ( 20 ; 30 ]. 3.3.2.8 Comentários sobre as respostas à Questão Discursiva 5 ( Bacharelado ) A questão 5 da prova do Bacharelado apresentava um problema simples, que é modelável por uma equação diofantina linear e cujas soluções são as soluções inteiras positivas da referida equação.
O conteúdo abordado na questão costuma ser trabalhado em disciplinas iniciais de Álgebra, que mais frequentemente se dão no primeiro ou segundo períodos dos cursos de Bacharelado.
No item ‘ a ’, foi solicitado ao estudante o apontamento de uma equação diofantina linear capaz de modelar o problema apresentado.
O padrão de resposta considerado para o item ‘ a ’ é a apresentação de uma das equações 9x + 15y = 246 ou 3x + 5y = 82, acompanhada da indicação de que x é o número de ingressos vendidos para estudantes, e y é o número de ingressos vendidos para o público geral.
O item ‘ b ’ solicitava ao estudante as soluções do problema exposto no enunciado, mas não exige que tais soluções fossem obtidas por meio da resolução formal da equação diofantina apontada no item ‘ a ’.
Essa é uma fragilidade da questão que impôs aos critérios de correção a necessidade de considerar a mera apresentação das soluções positivas, mesmo que de forma injustificada, presumidamente obtidas por testagens diretas junto à equação.
Originalmente, é claro, esperava -se que o estudante resolvesse a equação diofantina por meio das técnicas usuais para, ao final, escolher e indicar as suas soluções positivas.
O padrão de resposta considerado para o item ‘ b ’ consiste na resolução formal da equação diofantina apontada no item ‘ a ’, seguida da obtenção das cinco soluções positivas : ( x, y )  ( 24, 2 ), ( 19, 5 ), ( 14, 8 ), ( 9, 11 ), ( 4, 14 ).
O espaço destinado à resolução da questão foi adequado, e a dificuldade da questão pode ser considerada média.
O padrão de respostas da questão 5 era constituído por 3 etapas : 1 ) a indicação da equação diofantina adequada à modelagem do problema ( item ‘ a ’ ) ; 2 ) a resolução formal da referida equação diofantina e a obtenção da forma geral ( paramétrica ) de suas soluções ( item ‘ b ’ ) ; 3 ) a obtenção das soluções positivas a partir da forma geral ( paramétrica ) de suas soluções ( item ‘ b ’ ).
Como já dito anteriormente, o enunciado da questão não exigiu que as soluções positivas fossem obtidas a partir da resolução formal da equação diofantina.
Diante disso, houve estudantes que não cumpriram a etapa 2 e obtiveram as soluções positivas da equação por testagem, ou por outro modo informal.
Nesses casos, os estudantes que apresentaram as 5 soluções previstas no padrão de respostas receberam a pontuação integral no item ‘ b ’.
Alguns poucos estudantes utilizaram congruências modulares para encaminhar a solução ( corretamente ).
A média das notas obtidas nessa questão demonstrou que o desempenho foi mediano.
O fato de os estudantes que simplesmente apresentaram as 5 soluções positivas terem sido plenamente pontuados pode ter contribuído para o aumento dessa média.
De modo geral, as respostas consideradas fracas ( notas entre zero e 30 ) foram aquelas que apenas apresentaram a equação diofantina correta, acompanhada, ou não, da indicação correta do significado das variáveis x e y, solicitada na questão.
Respostas consideradas medianas ( notas entre 35 e 70 ) foram aquelas em que o estudante resolveu corretamente a equação diofantina ( correta, quase sempre ), mas não indicou as soluções positivas, ou as indicou em menor número.
As respostas boas ( notas entre 75 e 100 ) foram aquelas em que o estudante indicou corretamente a equação diofantina, resolveu -a formalmente e indicou soluções positivas da mesma.
De modo geral, os erros mais comumente cometidos pelos estudantes foram : 1 ) A indicação de uma equação diofantina inadequada.
Tal erro é, normalmente, decorrente da dificuldade de o estudante compreender o problema e selecionar suas variáveis de interesse.
Isso pode se dar por diversos motivos, dentre os quais se destacam a dificuldade de compreensão de textos básicos, assim como a falta de hábito na resolução de problemas ( estratégias de modelagem ). 2 ) Equívocos na resolução formal da equação diofantina obtida ( erros de conta, erros na identificação de uma solução particular, etc. ).
Erros técnicos resultantes de deficiências básicas na manipulação algébrica ou erros de conta. 3 ) A não indicação das soluções positivas, mesmo após a consideração da equação correta e de sua respectiva resolução formal ( correta ).
Esse tipo de erro pode ser decorrente da falta de atenção, ou da incompreensão de que as variáveis inteiras x e y poderiam assumir apenas valores positivos.
Esse tipo de erro corrobora a hipótese de que os estudantes não vêm trabalhando com a resolução de problemas em contextos mais gerais, mas, sim, apenas com a resolução de problemas em contexto puramente matemático.
Houve estudantes que indicaram uma equação diofantina errada, mas a resolveram formalmente de forma correta e indicaram, ao final, as suas soluções positivas.
Esses estudantes não receberam pontos no item ‘ a ’, mas receberam integralmente os pontos do item ‘ b ’.
Fato associado a definição de critérios de correção construídos de modo a não punir um estudante duas ou mais vezes, por conta de algum erro que se deu em um dos itens da questão e que se propagou aos itens seguintes.
De modo geral, a correção da questão confirmou a hipótese de que a formação dos Bacharéis em Matemática se dá em um perfil que privilegia mais a técnica de resolução das equações diofantinas do que outros pontos que a ela são anteriores ( modelagem ) e posteriores ( validação das soluções obtidas ).
Dentro desse perfil de formação, naturalmente, tornam -se aparentes apenas aqueles erros que decorrem das lacunas que os estudantes possuem nas manipulações elementares de ordem aritmética e algébrica.
A deficiência de formação nas técnicas de modelagem, que se coloca por meio de equívocos na escolha das variáveis e na adaptação e cumprimento das condições de contorno de um problema extra matemático, no entanto, foi responsável pela maior parte dos erros.
Tais deficiências passam despercebidas pelos 3.3.2.9 Análise da Questão Discursiva 5 do Componente de Conhecimento Específico ( Licenciatura ) A Tabela 3.21 contém as informações relativas à questão 5 do conjunto do Componente de Conhecimento Específico ( Licenciatura ).
O desempenho dos estudantes nessa questão foi superior ao das questões 3 e 4.
A nota média dos estudantes de todo o Brasil foi 14,2.
A maior média foi registrada na região Sul ( 15,8 ), enquanto a menor média foi registrada na região Norte ( 11,0 ).
Quanto à variabilidade das notas, o desvio padrão dos estudantes do Brasil, como um todo, foi 18,7.
Enquanto o maior desvio foi encontrado na região Nordeste ( 20,0 ), o menor foi encontrado na região Norte ( 15,6 ).
A nota máxima ( 100,0 ) e a nota mínima ( 0,0 ) são as mesmas em todas as regiões, com exceção da nota máxima da região Norte ( 90,0 ).
Para o conjunto de estudantes de Matemática ( Licenciatura ), a mediana foi 0,0, o que também ocorreu para as regiões Norte e Nordeste.
Para as demais regiões do Brasil a mediana foi 10,0.
O Gráfico 3.15 representa a distribuição das notas da questão discursiva 5 no Componente de Conhecimento Específico para estudantes ( Licenciatura ).
Mais uma vez, destaca -se o grande número de estudantes que deixaram a questão em branco, quase 40 \%, o valor modal. 3.3.2.10 Comentários sobre as respostas à Questão Discursiva 5 ( Licenciatura ) A questão propunha ao estudante a reflexão acerca de como a resolução de problemas se coloca e impacta as relações de ensino e aprendizagem da Matemática, mas de modo bastante específico e tematizado sobre um jogo bastante popular, a Torre de Hanói.
A questão era composta por três itens, que solicitavam : a ) o posicionamento de tal recurso nos \textbf{Parâmetros} Curriculares Nacionais, no que se refere ao ensino e à aprendizagem da Matemática ; b ) a indicação de conteúdos matemáticos que podem ser abordados por ele ; c ) a identificação da relação existente entre o número de discos e o número de movimentos mínimos a serem realizados no jogo para que o objetivo proposto seja alcançado.
Mais especificamente, o item ‘ a ’ convidava a apresentar objetivos que poderiam ser alcançados pelo professor, a partir da utilização da Torre de Hanoi, conforme é apontado pelos \textbf{Parâmetros} Curriculares Nacionais.
O item ‘ b ’ solicitava a indicação de três objetos matemáticos que poderiam ser abordados por meio do uso da Torre de Hanoi.
De certo modo, o item ‘ b ’ solicitava possibilidades mais relacionadas às perspectivas de ensino do que às perspectivas de aprendizagem.
Há, portanto, no que se refere ao eixo do ensino, uma \textbf{interseção} entre os temas tratados no item ‘ a ’ e no item ‘ b ’, uma vez que possíveis objetivos de utilização do recurso ( item ‘ a ’ ) são o desenvolvimento, a abordagem e a revisão de conceitos matemáticos específicos ( item ‘ b ’ ).
O item ‘ c ’ solicitava a determinação de uma fórmula para o número de movimentos mínimos ( com 3 pinos ), certamente conhecida/memorizada por muitos estudantes que já conheciam a Torre de Hanoi a priori.
A questão pode ser considerada adequada para estudantes que já conhecessem o material, Torre de Hanoi.
Aqueles que não a conheciam e contaram apenas com o enunciado da questão para conhecê -la, certamente não encontraram na prova o fórum adequado de manipulação do recurso, capaz de sustentar as inferências específicas acerca do potencial de seu uso, solicitadas nos itens ‘ a ’ e ‘ b ’.
Ou, até mesmo, capaz de esclarecer a mecânica do jogo a ponto de tornar viável a elaboração do argumento indutivo necessário à obtenção da fórmula solicitada no item ‘ c ’.
A questão poderia ser mais rica em termos \textbf{avaliativos}, se o recurso/jogo não tivesse sido fixado e se os dois primeiros itens solicitassem formalmente : ( a ) a reflexão sobre questões relativas ao ensino ; ( b ) a reflexão sobre questões relativas à aprendizagem ( ambos exigindo apontamentos dos PCN´s ).
O reconhecimento do potencial da Resolução de Problemas e da Modelagem no processo de ensino e da aprendizagem da Matemática é mais relevante para os futuros professores de Matemática do que as especificidades de um jogo ( Torre de Hanói ) em si.
Reafirma -se que, infelizmente, a questão privilegiou os estudantes que já conheciam o jogo antes da prova, em detrimento daqueles que não o conheciam, mesmo que os últimos tivessem uma visão interessante acerca da Resolução de Problemas, de modo mais geral.
O padrão de respostas era bastante amplo para o item ‘ a ’, uma vez que os PCN´s destacam os atributos da Resolução de Problemas em vieses de ensino e aprendizagem, a saber : a sua natural capacidade de introduzir, construir e revisitar conceitos matemáticos, de contextualizar práticas e objetos matemáticos, de construir significados, de motivar e estimular a aprendizagem colaborativa, etc.
O padrão de respostas do item ‘ b ’ considerava os seguintes conceitos matemáticos ( ou equivalentes ) : princípios de contagem, ordenação, conceito de função, função exponencial, função logarítmica, progressão geométrica, sequências e o trabalho com formas geométricas.
Este último não é um objeto matemático diretamente relacionado ao uso tradicional da Torre de Hanói, mas foi considerado pelo fato de muitos O padrão de resposta do item ‘ c ’ considerava a simples exposição da relação n(x) = 2^x − 1, uma vez que demonstrações por indução estariam fora de questão e os que não conheciam a Torre de Hanói conseguiriam, mais provavelmente, inferir tal fórmula a partir de uma tabela para x = 1, 2 e 3, ainda que a mesma não fosse suficiente para tal.
O espaço disponibilizado para responder à questão foi insuficiente, a se considerar a amplitude das soluções esperadas conforme previstas pelo padrão de resposta.
Em comparação com as demais questões de formação específica, a questão pode ser considerada difícil, pois exigia uma articulação de conhecimentos que ainda estão em construção, notoriamente, dentre a maioria dos estudantes.
As respostas dadas pelos estudantes variaram muitíssimo, conforme houvesse o conhecimento prévio da Torre de Hanói ou não.
A percepção de que os itens ‘ a ’ e ‘ b ’ possuíam \textbf{interseção} foi corroborada pelas inúmeras respostas que apresentaram conteúdos matemáticos já no item ‘ a ’.
Nesse ponto, a correção considerou as ideias apresentadas nos dois itens, pontuando elementos das respostas de um item que estivessem presentes na resposta do outro.
De modo geral, as respostas corretas dadas nos itens ‘ b ’ e ‘ c ’ seguiram o padrão de respostas considerado.
A média das notas dos alunos nessa questão foi baixíssima.
A média poderia ter sido ainda mais baixa se o critério de correção desconsiderasse respostas lacônicas e superficiais ou exigisse a diferenciação entre objetivos ( item ‘ a ’ ) e conteúdos ( item ‘ b ’ ).
De modo geral, as respostas consideradas fracas ( notas entre zero a 30 ) foram aquelas que se deram sobre um dos itens ‘ a ’ e ‘ b ’, de forma parcialmente correta.
Ficou claro que esta foi a categoria das respostas dadas pelos alunos que não conheciam a Torre de Hanoi.
As respostas consideradas medianas ( notas entre 35 e 70 ) foram aquelas em que os itens ‘ a ’ e ‘ b ’ foram abordados de forma mais correta.
As respostas consideradas boas ( notas entre 75 e 100 ) foram aquelas provenientes de estudantes que conheciam a Torre de Hanoi previamente e que conseguiram resgatar os pontos principais que foram solicitados.
Sobre o item ‘ a ’, de modo geral, os estudantes seguiram discursos já bastante populares sobre a utilização de jogos e atividades lúdicas, reproduzindo slogans, sem profundidade.
Muitos responderam por meio de frases únicas, como"Desenvolver o raciocínio lógico", sem demais considerações.
De toda forma, respostas lacônicas foram apenas parcialmente pontuadas. foram apenas parcialmente pontuadas.
A correção mostrou de forma notória a dificuldade de os estudantes escreverem textos simples, de modo objetivo : cometem inúmeros erros gramaticais, são redundantes, ambíguos e desconexos.
Os itens ‘ b ’ e ‘ c ’ foram respondidos de forma bastante diferente dentre os estudantes que já conheciam a Torre de Hanoi previamente.
Destaca -se, no entanto, a diferença entre as respostas dadas no item ‘ b ’ : os estudantes que não conheciam a Torre de Hanoi indicaram que o recurso seria útil para o desenvolvimento de noções sobre formas geométricas, área e perímetro de círculos, o trabalho com cores e atividades motoras.
Aqueles que já haviam manipulado a Torre de Hanoi sabiam que esses não eram os focos para o uso de tal material em situações de ensino.
Os alunos que não conheciam o recurso, a se julgar pela resposta dada no item ‘ b ’, não conseguiram responder ao item ‘ c ’, de modo geral.
Os alunos que já conheciam a Torre de Hanoi responderam, por vezes, apenas apresentando a fórmula N(𝑥) = 2𝑥−1, ou de outras similares, obtidas por equívocos na recuperação da fórmula correta da memória : N(𝑥) = 2𝑥−1, N(𝑥) = 2𝑥.
Alguns estudantes apresentaram uma tabela com o número de movimentos em função do número de discos, no entanto, é mais provável que eles a tenham apresentado para sustentar a fórmula que eles já sabiam de cor, em vez de, efetivamente, para inferi -la.
Pouquíssimos estudantes apresentaram a fórmula e a demonstraram ( por indução ).
De modo geral, as respostas encontradas revelaram que os estudantes tiveram dificuldades de encaminhar uma reflexão mais profunda acerca da resolução de problemas.
Mais precisamente, revelaram a falta de articulação existente entre os conhecimentos específicos ( em Matemática ) e os conhecimentos \textbf{pedagógicos} dos estudantes.
Essa falta de articulação é característica e persistente nos cursos de Licenciatura em Matemática há anos e decorre de propostas curriculares que ainda estão polarizadas entre o específico e o \textbf{pedagógico}, como ocorria nos antigos modelos 3+1 ( 3 anos de disciplinas da área específica + 1 ano de disciplinas \textbf{pedagógicas}, oferecidas de forma desvinculada da área específica ).
É fundamental que a formação do professor de Matemática reconheça não apenas a importância dos conhecimentos específicos e \textbf{pedagógicos}, mas, sobretudo, a articulação não trivial dos fundamentos da Matemática Escolar com especificidades filosóficas e \textbf{pedagógicas} essenciais, longe dos discursos fundamentalistas que advogam a favor de um ou de outro lado, por meio de slogans e discursos retóricos.
As respostas dadas no item ‘ a ’ são preocupantes, pois revelam não apenas a superficialidade demonstrada de forma lacônica, mas também a absoluta ausência de reflexões mais críticas.
Não houve respostas, por exemplo, que abordassem os eventuais limites da Torre de Hanoi.
De modo geral, as respostas se basearam mais sobre os seus pontos positivos referentes à aprendizagem do que sobre aqueles referentes ao ensino.
Por isso, é aparente que os estudantes do curso de Licenciatura em Matemática, futuros professores de Matemática, não articulam de modo mais profundo o recurso \textbf{pedagógico} com o conteúdo que será trabalhado com seus futuros alunos.

% matched lemmas: avaliativo, grosso, interseção, parâmetro, pedagógico
\end{textsample}
